
\documentclass[11pt,titlepage,a4paper]{report}

%INCLUSIONE PACCHETTI
%---------------------------------------------
\usepackage[italian]{babel}
\usepackage{fancyhdr}
\usepackage{graphicx}
\graphicspath{{./pics/}} % cartella di salvataggio immagini

% STILE DI PAGINA
%---------------------------------------------
\pagestyle{fancy}
\renewcommand{\sectionmark}[1]{\markright{\thesection.\ #1}}
\lhead{\nouppercase{\rightmark}}
\rhead{\nouppercase{\leftmark}}
\renewcommand{\chaptermark}[1]{%
\markboth{\thechapter.\ #1}{}}

%Ridefinisco lo stile plain della pagina
\fancypagestyle{plain}{%
	\lhead{\includegraphics[height=50pt]{logo.eps}}
	\chead{}
	\rhead{HappyCode inc \\ happycodeinc@gmail.com}
	\lfoot{BR-jsys}
	\cfoot{\thepage}
	\renewcommand{\headrulewidth}{1pt}
	\renewcommand{\footrulewidth}{1pt}
}
% layout
\begin{document}
%definizione variabili 
\newcommand{\lv}{ 0.3 } % latest version
\newcommand{\dt}{ Test Report }% Document title
%common variables
\newcommand{\br}{\underline{business rule}}
\newcommand{\brs}{\underline{business rules}}
\newcommand{\bo}{\underline{business object}}
\newcommand{\bos}{\underline{business objects}}
\newcommand{\rp}{\underline{repository}}
\newcommand{\brp}{BusinessRuleParser}
\newcommand{\brl}{BusinessRuleLexer}
\newcommand{\BR}{\underline{BusinessRule}}

%nomi dei componenti
\newcommand{\AT}{Alessia Trivellato}
\newcommand{\ET}{Elena Trivellato}
\newcommand{\FC}{Filippo Carraro}
\newcommand{\LA}{Luca Appon}
\newcommand{\MB}{Michele Bortolato}
\newcommand{\MT}{Marco Tessarotto}
\newcommand{\MM}{Mattia Meroi}%altre variabili
% ultime versioni dei documenti da modificare solo alla fine
\newcommand{\AR}{AnalisiDeiRequisiti.2.6.pdf}
\newcommand{\DdP}{DefinizioneDiProdotto.0.9.pdf}
\newcommand{\G}{ Glossario.1.8.pdf }
\newcommand{\NdP}{NormeDiProgetto.2.0.pdf}
\newcommand{\PdQ}{ PianoDiQualifica.2.0.pdf }
\newcommand{\PdP}{ PianoDiProgetto.1.7.pdf }
\newcommand{\ST}{SpecificaTecnica.1.5.pdf}
\newcommand{\TR}{TestReport.0.7.pdf}
\newcommand{\MU}{ManualeUtente.0.3.pdf}%nomi documenti
%fine definizione variabili
\hyphenation{
 a-na-lo-go
 as-so-cia-zio-ne
 %attività non si può inserire come tutte le parole accentate che vanno messe nel testo semplice scritte at\-ti\-vi\-tà o come variabile
 coe-ren-za
 com-po-nen-ti
 des-crit-te
 des-cri-zio-ni
 di-a-gram-ma
 di-a-gram-mi
 e-le-men-to
 e-se-gui-re
 e-si-sten-ti
 es-pli-ci-to
 glo-bal-men-te
 glos-sa-rio
 li-vel-lo
 ne-ces-sa-rio
 per-met-te-re
 re-po-si-to-ry
 re-vi-sio-na-men-to
 ri-chies-te
 se-gna-la-ta
 va-li-da-zio-ne
 va-ria-bi-li
 ve-ri-fi-ca-re
 vi-sua-liz-za-te
 e-ven-tua-li
 o-pe-ra-zio-ne
 ar-chi-via-zio-ne
 mo-di-fi-ca
}


%sillabazione
\begin{titlepage}
\begin{center}
\vspace*{0.5in}
\includegraphics{logo.eps}
\vspace*{0.2in}

{\Large \textbf{BR-jsys}}
{\Large \emph{business rules} per sistemi gestionali in architettura J2EE } 
\vspace{2in}

\Huge \textsc{ \dt }

\end{center}
\end{titlepage}
\vspace*{0.5in}%pagina del titolo


\begin{center}
\thispagestyle{plain}
\begin{table}[htbp]
\large{
\begin{tabular}{l}
\Large{\textbf{\textsf{Capitolato: ''BR-jsys``}}} \\
\begin{tabular}{|p{6cm}|p{6cm}|}
\hline
\textbf{Data creazione:} & 02/03/2008 \\ \hline
\textbf{Versione:} & \lv \\ \hline
\textbf{Stato del documento:} & Formale, esterno \\ \hline
% ----------------------------------------------------------------------------autori
\textbf{Revisione RQ} & \\ \hline
\textbf{Redazione:} & \MT \\ \hline
\textbf{Revisione:} &  \\ \hline
\textbf{Approvazione:}  &  \\ \hline
\end{tabular} \\
\end{tabular}
}
\end{table}

\begin{table}[hbtp]
\large{
\begin{tabular}{l}
\Large{\textbf{\textsf{Lista di distribuzione}}} \\
\begin{tabular}{|p{6cm}|p{6cm}|} \hline
%  -------------------------------------------------------------lista di distribuzione
{HappyCode inc}& Gruppo di lavoro \\ \hline
{Tullio Vardanega, Renato Conte}& Committenti \\ \hline 
{Zucchetti S.r.l}& Azienda proponente\\ \hline
\end{tabular} \\
\end{tabular}
}
\end{table}
\begin{table}[hbtp]

\Large{\textbf{\textsf{Diario delle modifiche}}} \\
\begin{small}
\begin{tabular}[t]{|p{1,2cm}|p{1.9cm}|p{2.9cm}|p{5cm}|} \hline
Versione & Data & Autore & Descrizione \\ \hline
%-------------------------------------------------------------------------------diario modifiche
0.7 & 07/03/2008 & \MT & Aggiunta dei 'test case' per le integrazioni di componenti.\\ \hline
0.6 & 07/03/2008 & \MM & Aggiunta la tabella di tracciamento 'requisiti - test'.\\ \hline
0.5 & 06/03/2008 & \LA & Aggiunto Sommario.\\ \hline
0.4 & 05/03/2008 & \MT & Aggiunta delle tabelle per contenere gli esiti, e dei 'test case' per il 'comunicatore' \\ \hline
0.3 & 04/03/2008 & \MT & Modifica alla nomenclatura dei test.\\ \hline
0.2 & 03/03/2008 & \MT & Aggiunta dei 'test case' per il 'validatore' .\\ \hline
0.1 & 02/03/2008 & \MT & Stesura iniziale del documento.\\ \hline
\end{tabular} \\
\end{small}


\end{table}
\end{center}
\newpage
\tableofcontents

\chapter*{Sommario}
\addcontentsline{toc}{chapter}{ I Sommario}
Il presente documento illustra le strategie utilizzate per testare il prodotto BR-jsys in ogni sua componente.

\chapter*{Glossario}
\addcontentsline{toc}{chapter}{II Glossario}
Per gli altri termini non citati nella tabella soprastante si fa riferimento al file esterno \G. I termini presenti in questo documento la cui descrizione \`e riportata nel glossario sono evidenziati mediante sottolineatura.

\chapter{Introduzione}

\section{Scopo del documento}
Il presente documento viene redatto al fine di illustrare la progettazione dei \textit{test case} nel dettaglio indicando un insieme di input e il corrispondente insieme di \textit{output attesi}. Verranno inoltre riportati nel presente documento gli esiti dei test eseguiti sul sistema ``Br-jsys'' al fine di non appesantire la lettura del documento \PdQ .

\section{Definizioni, acronimi, abbreviazioni}
Nella tabella di seguito vengono riportate tutte le abbreviazioni utlizzate nel documento, comprese le sigle utilizzate per l'identificazione dei test.
\begin{center}
\begin{tabular}{||p{3.0cm}||p{8.5cm}||} \hline
\textbf{Abbreviazione} & \textbf{Significato} \\ \hline

VT\textit{\#} & Con VT seguito da un numero intero sono indicati i test che riguardano la 'macrocomponente validatore'.\\ \hline
CT\textit{\#} & Con CT seguito da un numero intero sono indicati i test che riguardano la 'macrocomponente comunicatore'.\\ \hline
GT\textit{\#} & Con GT seguito da un numero intero sono indicati i test che riguardano la 'macrocomponente gui'.\\ \hline
CVT\textit{\#} & Con CVT seguito da un numero intero sono indicati i test di integrazione tra 'validatore' e 'comunicatore'.\\ \hline
GVT\textit{\#} & Con GVT seguito da un numero intero sono indicati i test integrazione tra 'gui' e 'validatore'.\\ \hline
S\textit{\#} & Con S seguito da un numero intero sono indicati i test di integrazione del sistema.\\ \hline
RT\textit{} & Con RT si intende la riesecuzione di un test (test di regressione) avvenuta a fronte di una modifica al codice.\\ \hline

\end{tabular} \\
\end{center}

\chapter{Descrizione dei test case}
Vengono di seguito presentati i \textit{test case}. Ogni test case \`e composto da un input, un output atteso. Ogni test case \`e inoltre accompagnato da un breve testo che identifica lo scopo del test, e dall'elenco dei requisiti soddisfatti dal test.

\section{Test sulle componenti}
Vengono di seguito riportati i test riguardanti le macrocomponenti del sistema \underline{br-jsys}. Questi test si occupano di testare il corretto funzionamento della componente, nonch\`e di mostrare il soddisfacimento di alcuni requisiti.

\subsection{Test sul validatore}
La macrocomponente validatore verr\`a testata a scatola aperta ( whitebox ), inserendo in input al validatore un oggetto \textit{\br\ } con gli attributi di seguito specificati e confrontando l'output del validatore con l'output atteso.\\
\\
\begin{Large}\textbf{VT1a}\end{Large} \\
Si testa una \br\ semplice per verificare che siano accettati tutti i campi dati.\\
Requisiti soddisfatti dall'esecuzione: F1,F2,F10,F11,NU3
\begin{center}
\begin{tabular}{|p{5cm}|p{6cm}|} \hline
\textbf{Attributo \br} & \textbf{Valore} \\ \hline
Nome & BR\\ \hline
Commento & UnCommento\\ \hline
Business object associato & Articolo\\ \hline
Regola & entrate=0 AND NOT(scontato)=true\\ \hline
\end{tabular} \\
\end{center}
\begin{center}
\begin{tabular}{|p{11cm}|} \hline
\textbf{Output atteso}\\ \hline
\textless BusinessRule associated=``Articolo'' comment=``Commento'' name=``BR'' rule=``entrate=0 AND NOT(scontato)=true''\textgreater \\
\textless AstRuleVersion value=``(AND (= entrate 0) (= (NOT scontato) true))''/\textgreater \\
 \textless OpRule type=``AND''\textgreater \\
 \textless Conf type=``=''\textgreater \\
 \textless FIELD\textgreater entrate\textless /FIELD\textgreater \\
 \textless FLOAT\textgreater 0\textless /FLOAT\textgreater \\
 \textless /Conf\textgreater \\
 \textless Conf type=``=''\textgreater \\
\textless BoFun type=``NOT''\textgreater \\
 \textless FIELD\textgreater scontato\textless /FIELD\textgreater \\
\textless /BoFun\textgreater \\
 \textless BOOL\textgreater true\textless /BOOL\textgreater \\
\textless /Conf\textgreater \\
\textless /OpRule\textgreater \\
\textless /BusinessRule\textgreater \\
 \hline
\end{tabular} \\
\end{center}

\begin{Large}\textbf{VT1b}\end{Large} \\
Si testa una \br\ semplice per verificare che il commento pu\`o non essere inserito.\\
Requisiti soddisfatti dall'esecuzione: F1,F2,F10
\begin{center}
\begin{tabular}{|p{5cm}|p{6cm}|} \hline
\textbf{Attributo \br} & \textbf{Valore} \\ \hline
Nome & BR\\ \hline
Commento & \\ \hline
Business object associato & Articolo\\ \hline
Regola & entrate=3 AND uscite=4\\ \hline
\end{tabular} \\
\end{center}
\begin{center}
\begin{tabular}{|p{11cm}|} \hline
\textbf{Output atteso}\\ \hline
\textless BusinessRule associated=``Articolo'' comment=``'' name=``BR'' rule=``entrate=3 AND uscite=4''\textgreater \\
 \textless AstRuleVersion value=``(AND (= entrate 3) (= uscite 4))''/\textgreater \\
\textless OpRule type=``AND''\textgreater \\
\textless Conf type=``=``\textgreater \\
\textless FIELD\textgreater entrate\textless /FIELD\textgreater \\
\textless FLOAT\textgreater 3\textless /FLOAT\textgreater \\
\textless /Conf\textgreater \\
\textless Conf type=``=``\textgreater \\
\textless FIELD\textgreater uscite\textless /FIELD\textgreater \\
\textless FLOAT\textgreater 4\textless /FLOAT\textgreater \\
\textless /Conf\textgreater \\
\textless /OpRule\textgreater \\
\textless /BusinessRule\textgreater \\
 \hline
\end{tabular} \\
\end{center}

\begin{Large}\textbf{VT1c}\end{Large} \\
Si testa una \br\ semplice per verificare che il nome non \`e opzionale.\\
Requisiti soddisfatti dall'esecuzione: F1,F2,F10
\begin{center}
\begin{tabular}{|p{5cm}|p{6cm}|} \hline
\textbf{Attributo \br} & \textbf{Valore} \\ \hline
Nome & \\ \hline
Commento & Commento\\ \hline
Business object associato & Articolo\\ \hline
Regola & entrate=5 AND uscite=3\\ \hline
\end{tabular} \\
\end{center}
\begin{center}
\begin{tabular}{|p{11cm}|} \hline
\textbf{Output atteso}\\ \hline
\textless BusinessRule associated=``Articolo'' comment=``Commento'' name=``'' rule=``entrate=5 AND uscite=3''\textgreater \\
\textless AstRuleVersion value=``(AND (= entrate 5) (= uscite 3))''/\textgreater \\
\textless OpRule type=``AND''\textgreater \\
\textless Conf type=``=``\textgreater \\
\textless FIELD\textgreater entrate\textless /FIELD\textgreater \\
\textless FLOAT\textgreater 5\textless /FLOAT\textgreater \\
\textless /Conf\textgreater\\
 \textless Conf type=``=``\textgreater \\
\textless FIELD\textgreater uscite\textless /FIELD\textgreater \\
\textless FLOAT\textgreater 3\textless /FLOAT\textgreater \\
\textless /Conf\textgreater \\
\textless /OpRule\textgreater \\
\textless /BusinessRule\textgreater \\
 \hline
\end{tabular} \\
\end{center}

\begin{Large}\textbf{VT1d}\end{Large} \\
Si testa una \br\ semplice per verificare che il \bo\ non \`e opzionale.\\
Requisiti soddisfatti dall'esecuzione: F1,F2,F10,F11
\begin{center}
\begin{tabular}{|p{5cm}|p{6cm}|} \hline
\textbf{Attributo \br} & \textbf{Valore} \\ \hline
Nome & BR \\ \hline
Commento & Commento\\ \hline
Business object associato & \\ \hline
Regola & entrate=3 OR uscite=5\\ \hline
\end{tabular} \\
\end{center}
\begin{center}
\begin{tabular}{|p{11cm}|} \hline
\textbf{Output atteso}\\ \hline
\textless BusinessRule associated=``Articolo'' comment=``Commento'' name=``BR'' rule=``entrate=3 OR uscite=5''\textgreater\\
 \textless AstRuleVersion value=``(OR (= entrate 3) (= uscite 5))''/\textgreater\\
 \textless OpRule type=``OR''\textgreater \\
\textless Conf type=``=``\textgreater\\
 \textless FIELD\textgreater entrate\textless /FIELD\textgreater\\
 \textless FLOAT\textgreater 3\textless /FLOAT\textgreater\\
 \textless /Conf\textgreater\\
 \textless Conf type=``=``\textgreater\\
 \textless FIELD\textgreater uscite\textless /FIELD\textgreater\\
 \textless FLOAT\textgreater 5\textless /FLOAT\textgreater\\
 \textless /Conf\textgreater\\
 \textless /OpRule\textgreater\\
 \textless /BusinessRule\textgreater \\
 \hline
\end{tabular} \\
\end{center}

\begin{Large}\textbf{VT1e}\end{Large} \\
Si testa una \br\ semplice per verificare che la regola stessa non \`e opzionale.\\
Requisiti soddisfatti dall'esecuzione: F1,F2
\begin{center}
\begin{tabular}{|p{5cm}|p{6cm}|} \hline
\textbf{Attributo \br} & \textbf{Valore} \\ \hline
Nome & BR \\ \hline
Commento & Commento\\ \hline
Business object associato & Articolo \\ \hline
Regola & \\ \hline
\end{tabular} \\
\end{center}
\begin{center}
\begin{tabular}{|p{11cm}|} \hline
\textbf{Output atteso}\\ \hline
Errore sintattico\\
 \hline
\end{tabular} \\
\end{center}

\begin{Large}\textbf{VT1f}\end{Large} \\
Si testa una \br\ semplice per verificare che non vengano accettate \br\ vuote.\\
Requisiti soddisfatti dall'esecuzione: F1,F2
\begin{center}
\begin{tabular}{|p{5cm}|p{6cm}|} \hline
\textbf{Attributo \br} & \textbf{Valore} \\ \hline
Nome &  \\ \hline
Commento & \\ \hline
Business object associato &  \\ \hline
Regola & \\ \hline
\end{tabular} \\
\end{center}
\begin{center}
\begin{tabular}{|p{11cm}|} \hline
\textbf{Output atteso}\\ \hline
Errore sintattico\\
 \hline
\end{tabular} \\
\end{center}

\begin{Large}\textbf{VT2a}\end{Large} \\
Si testa una \br\ per verificare che vengano ben interpretate le parentesi tra 'AND' e 'OR'.\\
Requisiti soddisfatti dall'esecuzione: F1,F2,F10
\begin{center}
\begin{tabular}{|p{5cm}|p{6cm}|} \hline
\textbf{Attributo \br} & \textbf{Valore} \\ \hline
Nome & BR \\ \hline
Commento & Commento\\ \hline
Business object associato & Articolo \\ \hline
Regola & (uscite=5 OR uscite=5)  AND ( uscite=5 OR uscite=5)  \\ \hline
\end{tabular} \\
\end{center}
\begin{center}
\begin{tabular}{|p{11cm}|} \hline
\textbf{Output atteso}\\ \hline
\textless BusinessRule associated=``Articolo'' comment=``Commento'' name=``BR'' rule=``(uscite=5 OR uscite=5)  AND ( uscite=5 OR uscite=5)''\textgreater\\
 \textless AstRuleVersion value=``(AND (OR (= uscite 5) (= uscite 5)) (OR (= uscite 5) (= uscite 5)))''/\textgreater \\
\textless OBool type=``AND''\textgreater\\
\textless OBool type=``OR''\textgreater\\
\textless Conf type=``=``\textgreater \\
\textless FIELD\textgreater uscite\textless /FIELD\textgreater\\
 \textless FLOAT\textgreater 5\textless /FLOAT\textgreater\\
 \textless /Conf\textgreater\\
\textless Conf type=``=``\textgreater \\
\textless FIELD\textgreater uscite\textless /FIELD\textgreater\\
 \textless FLOAT\textgreater 5\textless /FLOAT\textgreater\\
 \textless /Conf\textgreater\\
\textless /OBool\textgreater \\
\textless OBool type=``OR''\textgreater\\
\textless Conf type=``=``\textgreater \\
\textless FIELD\textgreater uscite\textless /FIELD\textgreater\\
 \textless FLOAT\textgreater 5\textless /FLOAT\textgreater\\
 \textless /Conf\textgreater\\
\textless Conf type=``=``\textgreater \\
\textless FIELD\textgreater uscite\textless /FIELD\textgreater\\
 \textless FLOAT\textgreater 5\textless /FLOAT\textgreater\\
 \textless /Conf\textgreater\\
\textless /OBool\textgreater \\
\textless /OBool\textgreater \\
\textless /BusinessRule\textgreater \\
 \hline
\end{tabular} \\
\end{center}

\begin{Large}\textbf{VT2b}\end{Large} \\
Si testa una \br\ per verificare che vengano corretamente riportato il susseguirsi di 'AND' e 'OR'.\\
Requisiti soddisfatti dall'esecuzione: F1,F2,F10
\begin{center}
\begin{tabular}{|p{5cm}|p{6cm}|} \hline
\textbf{Attributo \br} & \textbf{Valore} \\ \hline
Nome & BR \\ \hline
Commento & Commento\\ \hline
Business object associato & Articolo \\ \hline
Regola & uscite=5 AND uscite=5 OR uscite=5 AND uscite=5 OR uscite=5 AND uscite=5 OR uscite=5 AND uscite=5 AND uscite=5 AND uscite=5 OR uscite=5 OR uscite=5 AND uscite=5 OR uscite=5 OR uscite=5 AND uscite=5 \\ \hline
\end{tabular} \\
\end{center}
\begin{center}
\begin{tabular}{|p{11cm}|} \hline
\textbf{Output atteso}\\ \hline
\textless BusinessRule associated=``Articolo'' comment=``Commento'' name=``BR'' rule=``uscite=5 OR uscite=5  AND  uscite=5 OR uscite=5''\textgreater\\
 \textless AstRuleVersion value=``(OR (AND (OR (= uscite 5) (= uscite 5)) (= uscite 5)) (= uscite 5))''/\textgreater \\
\textless OBool type=``OR''\textgreater\\
\textless OBool type=``AND''\textgreater\\
\textless OBool type=``OR''\textgreater\\
\textless Conf type=``=``\textgreater \\
\textless FIELD\textgreater uscite\textless /FIELD\textgreater\\
 \textless FLOAT\textgreater 5\textless /FLOAT\textgreater\\
 \textless /Conf\textgreater\\
\textless Conf type=``=``\textgreater \\
\textless FIELD\textgreater uscite\textless /FIELD\textgreater\\
 \textless FLOAT\textgreater 5\textless /FLOAT\textgreater\\
 \textless /Conf\textgreater\\
\textless /OBool\textgreater \\
\textless Conf type=``=``\textgreater \\
\textless FIELD\textgreater uscite\textless /FIELD\textgreater\\
 \textless FLOAT\textgreater 5\textless /FLOAT\textgreater\\
 \textless /Conf\textgreater\\
\textless /OBool\textgreater \\
\textless Conf type=``=``\textgreater \\
\textless FIELD\textgreater uscite\textless /FIELD\textgreater\\
 \textless FLOAT\textgreater 5\textless /FLOAT\textgreater\\
 \textless /Conf\textgreater\\
\textless /OBool\textgreater \\
\textless /BusinessRule\textgreater \\
 \hline
\end{tabular} \\
\end{center}

\begin{Large}\textbf{VT2c}\end{Large} \\
Si testa una \br\ per verificare che possano essere assenti 'AND' e 'OR'.\\
Requisiti soddisfatti dall'esecuzione: F1,F2,F10
\begin{center}
\begin{tabular}{|p{5cm}|p{6cm}|} \hline
\textbf{Attributo \br} & \textbf{Valore} \\ \hline
Nome & BR \\ \hline
Commento & Commento\\ \hline
Business object associato & Articolo \\ \hline
Regola & entrate=5+2\\ \hline
\end{tabular} \\
\end{center}
\begin{center}
\begin{tabular}{|p{11cm}|} \hline
\textbf{Output atteso}\\ \hline
\textless BusinessRule associated=``Articolo'' comment=``Commento'' name=``BR'' rule=``entrate=5+2''\textgreater \\
\textless AstRuleVersion value=``(= entrate (+ 5 2))''/\textgreater \\
\textless Bconf type=``=``\textgreater\\
 \textless FIELD\textgreater entrate\textless /FIELD\textgreater \\
\textless OFloat type=``+''\textgreater \\
\textless FLOAT\textgreater 5\textless /FLOAT\textgreater \\
\textless FLOAT\textgreater 2\textless /FLOAT\textgreater\\
 \textless /OFloat\textgreater \\
\textless /Bconf\textgreater \\
\textless /BusinessRule\textgreater \\
 \hline
\end{tabular} \\
\end{center}

\begin{Large}\textbf{VT3a}\end{Large} \\
Si testa una \br\ per verificare che il 'message' sia opzionale.\\
Requisiti soddisfatti dall'esecuzione: F1,F2,F10
\begin{center}
\begin{tabular}{|p{5cm}|p{6cm}|} \hline
\textbf{Attributo \br} & \textbf{Valore} \\ \hline
Nome & BR \\ \hline
Commento & Commento\\ \hline
Business object associato & Articolo \\ \hline
Regola & entrate=SUM(prezzoBase)\\ \hline
\end{tabular} \\
\end{center}
\begin{center}
\begin{tabular}{|p{11cm}|} \hline
\textbf{Output atteso}\\ \hline
\textless BusinessRule associated=``Articolo'' comment=``Commento'' name=``BR'' rule=``entrate=SUM(prezzoBase)''\textgreater \\
\textless AstRuleVersion value=``(= entrate (SUM prezzoBase))''/\textgreater \\
\textless Bconf type=``=``\textgreater \\
\textless FIELD\textgreater entrate\textless /FIELD\textgreater \\
\textless FlFun type=``SUM''\textgreater \\
\textless FIELD\textgreater prezzoBase\textless /FIELD\textgreater \\
\textless /FlFun\textgreater \\
\textless /Bconf\textgreater \\
\textless /BusinessRule\textgreater \\
 \hline
\end{tabular} \\
\end{center}

\begin{Large}\textbf{VT3b}\end{Large} \\
Si testa una \br\ per verificare che sia possibile inserire un 'message'.\\
Requisiti soddisfatti dall'esecuzione: F1,F2,F10,NU3
\begin{center}
\begin{tabular}{|p{5cm}|p{6cm}|} \hline
\textbf{Attributo \br} & \textbf{Valore} \\ \hline
Nome & BR \\ \hline
Commento & Commento\\ \hline
Business object associato & Articolo \\ \hline
Regola & entrate=AVG(prezzoBase)+2 message(``errore nel confronto'')\\ \hline
\end{tabular} \\
\end{center}
\begin{center}
\begin{tabular}{|p{11cm}|} \hline
\textbf{Output atteso}\\ \hline
\textless BusinessRule associated=``Articolo'' comment=``Commento'' name=``BR'' rule=``entrate=AVG(prezzoBase)+2 message(\&quot;errore nel confronto\&quot;)''\textgreater \\
\textless AstRuleVersion value=``(= entrate (+ (AVG prezzoBase) 2) (message \&quot;errore nel confronto\&quot;))''/\textgreater\\
 \textless Bconf type=``=``\textgreater \\
\textless FIELD\textgreater entrate\textless /FIELD\textgreater \\
\textless OFloat type=``+''\textgreater \\
\textless FlFun type=``AVG''\textgreater \\
\textless FIELD\textgreater prezzoBase\textless /FIELD\textgreater \\
\textless /FlFun\textgreater \\
\textless FLOAT\textgreater 2\textless /FLOAT\textgreater \\
\textless /OFloat\textgreater \\
\textless Msg type=``message''\textgreater \\
\textless STRING\textgreater ``errore nel confronto''\textless /STRING\textgreater\\
 \textless /Msg\textgreater\\
 \textless /Bconf\textgreater \\
\textless /BusinessRule\textgreater \\
 \hline
\end{tabular} \\
\end{center}

\begin{Large}\textbf{VT4}\end{Large} \\
Si testa una \br\ per verificare che sia possibile fare confronti del tipo '\textless', '\textgreater', '\textless =', '\textgreater =', '=', '!='.\\
Requisiti soddisfatti dall'esecuzione: F1,F2,F10
\begin{center}
\begin{tabular}{|p{5cm}|p{6cm}|} \hline
\textbf{Attributo \br} & \textbf{Valore} \\ \hline
Nome & BR \\ \hline
Commento & Commento\\ \hline
Business object associato & Articolo \\ \hline
Regola & uscite\textless2 AND uscite\textgreater 2 AND uscite\textless=2 AND uscite\textgreater =2 AND uscite=2 AND uscite!=2\\ \hline
\end{tabular} \\
\end{center}
\begin{center}
\begin{tabular}{|p{11cm}|} \hline
\textbf{Output atteso}\\ \hline
\textless BusinessRule associated=``Articolo'' comment=``Commento'' name=``BR'' rule=``uscite\&lt;2 AND uscite\&gt; 2 AND uscite\&lt;=2 AND uscite\&gt;=2 AND uscite=2 AND uscite!=2''\textgreater\\
 \textless AstRuleVersion value=``(AND (AND (AND (AND (AND (= uscite 2) (!= uscite 2)) (\&gt;= uscite 2)) (\&lt;= uscite 2)) (\&gt; uscite 2)) (\&lt; uscite 2))''/\textgreater \\
\textless OBool type=``AND''\textgreater \textless OBool type=``AND''\textgreater\\
\textless OBool type=``AND''\textgreater \textless OBool type=``AND''\textgreater\\
\textless OBool type=``AND''\textgreater\\
\textless Conf type=``=``\textgreater \\
\textless FIELD\textgreater uscite\textless /FIELD\textgreater\\
 \textless FLOAT\textgreater 5\textless /FLOAT\textgreater\\
 \textless /Conf\textgreater\\
\textless Conf type=``!=``\textgreater \\
\textless FIELD\textgreater uscite\textless /FIELD\textgreater\\
 \textless FLOAT\textgreater 5\textless /FLOAT\textgreater\\
 \textless /Conf\textgreater\\
\textless /OBool\textgreater \\
\textless Conf type=``\&gt;=``\textgreater \\
\textless FIELD\textgreater uscite\textless /FIELD\textgreater\\
 \textless FLOAT\textgreater 5\textless /FLOAT\textgreater\\
 \textless /Conf\textgreater\\
\textless /OBool\textgreater \\
\textless Conf type=``\&lt;=``\textgreater \\
\textless FIELD\textgreater uscite\textless /FIELD\textgreater\\
 \textless FLOAT\textgreater 5\textless /FLOAT\textgreater\\
 \textless /Conf\textgreater\\
\textless /OBool\textgreater \\
\textless Conf type=``\&gt;``\textgreater \\
\textless FIELD\textgreater uscite\textless /FIELD\textgreater\\
 \textless FLOAT\textgreater 5\textless /FLOAT\textgreater\\
 \textless /Conf\textgreater\\
\textless /OBool\textgreater \\
\textless Conf type=``\&lt;``\textgreater \\
\textless FIELD\textgreater uscite\textless /FIELD\textgreater\\
 \textless FLOAT\textgreater 5\textless /FLOAT\textgreater\\
 \textless /Conf\textgreater\\
\textless /OBool\textgreater \\
\textless /BusinessRule\textgreater \\
 \hline
\end{tabular} \\
\end{center}


\begin{Large}\textbf{VT5a}\end{Large} \\
Si testa una \br\ per verificare che sia possibile fare operazioni aritmetiche semplici tra scalari.\\
Requisiti soddisfatti dall'esecuzione: F1,F2,F10
\begin{center}
\begin{tabular}{|p{5cm}|p{6cm}|} \hline
\textbf{Attributo \br} & \textbf{Valore} \\ \hline
Nome & BR \\ \hline
Commento & Commento\\ \hline
Business object associato & Articolo \\ \hline
Regola & uscite=(2+2-(3*4))/5 \\ \hline
\end{tabular} \\
\end{center}
\begin{center}
\begin{tabular}{|p{11cm}|} \hline
\textbf{Output atteso}\\ \hline
\textless BusinessRule associated=``Articolo'' comment=``Commento'' name=``BR'' rule=``uscite=(2+2-(3*4))/5''\textgreater \\
\textless AstRuleVersion value=``(= uscite (/ (- (+ 2 2) (* 3 4)) 5))''/\textgreater \\
 \textless Bconf type=``=``\textgreater \\
\textless FIELD\textgreater uscite\textless /FIELD\textgreater \\
\textless OFloat type=``/''\textgreater \\
\textless OFloat type=``-''\textgreater \\
\textless OFloat type=``+''\textgreater \\
\textless FLOAT\textgreater 2\textless /FLOAT\textgreater \\
\textless FLOAT\textgreater 2\textless /FLOAT\textgreater \\
\textless /OFloat\textgreater \\
\textless OFloat type=``*''\textgreater \\
\textless FLOAT\textgreater 3\textless /FLOAT\textgreater\\
\textless FLOAT\textgreater 4\textless /FLOAT\textgreater \\
\textless /OFloat\textgreater \\
\textless /OFloat\textgreater \\
\textless FLOAT\textgreater 5\textless /FLOAT\textgreater \\
\textless /OFloat\textgreater \\
\textless /Bconf\textgreater \\
\textless /BusinessRule\textgreater \\
 \hline
\end{tabular} \\
\end{center}

\begin{Large}\textbf{VT5b}\end{Large} \\
Si testa una \br\ per verificare che sia possibile fare operazioni aritmetiche semplici tra scalari e campi dati scalari.\\
Requisiti soddisfatti dall'esecuzione: F1,F2,F10
\begin{center}
\begin{tabular}{|p{5cm}|p{6cm}|} \hline
\textbf{Attributo \br} & \textbf{Valore} \\ \hline
Nome & BR \\ \hline
Commento & Commento\\ \hline
Business object associato & Articolo \\ \hline
Regola & uscite=(2+entrate-(uscite*4))/entrate \\ \hline
\end{tabular} \\
\end{center}
\begin{center}
\begin{tabular}{|p{11cm}|} \hline
\textbf{Output atteso}\\ \hline
\textless BusinessRule associated=``Articolo'' comment=``Commento'' name=``BR'' rule=``uscite=(2+entrate-(uscite*4))/entrate''\textgreater\\
 \textless AstRuleVersion value=``(= uscite (/ (- (+ 2 entrate) (* uscite 4)) entrate))''/\textgreater\\
 \textless Bconf type=``=``\textgreater \\
\textless FIELD\textgreater uscite\textless /FIELD\textgreater \\
\textless OFloat type=``/''\textgreater \\
\textless OFloat type=``-''\textgreater \\
\textless OFloat type=``+''\textgreater \\
\textless FLOAT\textgreater 2\textless /FLOAT\textgreater \\
\textless FIELD\textgreater entrate\textless /FIELD\textgreater \\
\textless /OFloat\textgreater \\
\textless OFloat type=``*''\textgreater \\
\textless FIELD\textgreater uscite\textless /FIELD\textgreater \\
\textless FLOAT\textgreater 4\textless /FLOAT\textgreater \\
\textless /OFloat\textgreater\\
 \textless /OFloat\textgreater \\
\textless FIELD\textgreater entrate\textless /FIELD\textgreater \\
\textless /OFloat\textgreater\\
 \textless /Bconf\textgreater \\
\textless /BusinessRule\textgreater \\
 \hline
\end{tabular} \\
\end{center}

\begin{Large}\textbf{VT5c}\end{Large} \\
Si testa una \br\ per verificare che sia possibile fare operazioni aritmetiche semplici tra campi dati scalari.\\
Requisiti soddisfatti dall'esecuzione: F1,F2,F10
\begin{center}
\begin{tabular}{|p{5cm}|p{6cm}|} \hline
\textbf{Attributo \br} & \textbf{Valore} \\ \hline
Nome & BR \\ \hline
Commento & Commento\\ \hline
Business object associato & Articolo \\ \hline
Regola & uscite=(uscite+entrate-(uscite*entrate))/entrate \\ \hline
\end{tabular} \\
\end{center}
\begin{center}
\begin{tabular}{|p{11cm}|} \hline
\textbf{Output atteso}\\ \hline
\textless BusinessRule associated=``Articolo'' comment=``Commento'' name=``BR'' rule=``uscite=(uscite+entrate-(uscite*entrate))/entrate''\textgreater\\
 \textless AstRuleVersion value=``(= uscite (/ (- (+ uscite entrate) (* uscite entrate)) entrate))''/\textgreater \\
\textless Bconf type=``=``\textgreater \\
\textless FIELD\textgreater uscite\textless /FIELD\textgreater\\
 \textless OFloat type=``/''\textgreater \\
\textless OFloat type=``-''\textgreater \\
\textless OFloat type=``+''\textgreater \\
\textless FIELD\textgreater uscite\textless /FIELD\textgreater\\
 \textless FIELD\textgreater entrate\textless /FIELD\textgreater \\
\textless /OFloat\textgreater \\
\textless OFloat type=``*''\textgreater \\
\textless FIELD\textgreater uscite\textless /FIELD\textgreater \\
\textless FIELD\textgreater entrate\textless /FIELD\textgreater \\
\textless /OFloat\textgreater \\
\textless /OFloat\textgreater \\
\textless FIELD\textgreater entrate\textless /FIELD\textgreater \\
\textless /OFloat\textgreater \\
\textless /Bconf\textgreater \\
\textless /BusinessRule\textgreater \\
 \hline
\end{tabular} \\
\end{center}

\begin{Large}\textbf{VT6a}\end{Large} \\
Si testa una \br\ per verificare che sia possibile fare la operazioni aritmetiche semplici tra scalare e vettore.\\
Requisiti soddisfatti dall'esecuzione: F1,F2,F10,NU1
\begin{center}
\begin{tabular}{|p{5cm}|p{6cm}|} \hline
\textbf{Attributo \br} & \textbf{Valore} \\ \hline
Nome & BR \\ \hline
Commento & Commento\\ \hline
Business object associato & Articolo \\ \hline
Regola & uscite=(2+prezzoBase) +(prezzoBase*4) +(prezzoBase/4) +(prezzoBase-3) \\ \hline
\end{tabular} \\
\end{center}
\begin{center}
\begin{tabular}{|p{11cm}|} \hline
\textbf{Output atteso}\\ \hline
\textless BusinessRule associated=``Articolo'' comment=``Commento'' name=``BR'' rule=``uscite=(2+prezzoBase)+(prezzoBase*4) +(prezzoBase/4)+(prezzoBase-3)''\textgreater\\
 \textless AstRuleVersion value=``(= uscite (+ (+ (+ (+ 2 prezzoBase) (* prezzoBase 4)) (/ prezzoBase 4)) (- prezzoBase 3)))''/\textgreater \\
\textless Bconf type=``=``\textgreater \\
\textless FIELD\textgreater uscite\textless /FIELD\textgreater \\
\textless OFloat type=``+''\textgreater \\
\textless OFloat type=``+''\textgreater \\
\textless OFloat type=``+''\textgreater \\
\textless OFloat type=``+''\textgreater \\
\textless FLOAT\textgreater 2\textless /FLOAT\textgreater \\
\textless FIELD\textgreater prezzoBase\textless /FIELD\textgreater \\
\textless /OFloat\textgreater \\
\textless OFloat type=``*''\textgreater \\
\textless FIELD\textgreater prezzoBase\textless /FIELD\textgreater \\
\textless FLOAT\textgreater 4\textless /FLOAT\textgreater \\
\textless /OFloat\textgreater \\
\textless /OFloat\textgreater \\
\textless OFloat type=``/''\textgreater\\
 \textless FIELD\textgreater prezzoBase\textless /FIELD\textgreater \\
\textless FLOAT\textgreater 4\textless /FLOAT\textgreater \\
\textless /OFloat\textgreater \\
\textless /OFloat\textgreater \\
\textless OFloat type=``-''\textgreater \\
\textless FIELD\textgreater prezzoBase\textless /FIELD\textgreater \\
\textless FLOAT\textgreater 3\textless /FLOAT\textgreater \\
\textless /OFloat\textgreater\\
 \textless /OFloat\textgreater\\
 \textless /Bconf\textgreater \\
\textless /BusinessRule\textgreater \\
 \hline
\end{tabular} \\
\end{center}

\begin{Large}\textbf{VT6b}\end{Large} \\
Si testa una \br\ per verificare che sia possibile fare la operazioni aritmetiche semplici tra campi dati scalari e vettori.\\
Requisiti soddisfatti dall'esecuzione: F1,F2,F10,NU1
\begin{center}
\begin{tabular}{|p{5cm}|p{6cm}|} \hline
\textbf{Attributo \br} & \textbf{Valore} \\ \hline
Nome & BR \\ \hline
Commento & Commento\\ \hline
Business object associato & Articolo \\ \hline
Regola & (entrate+prezzoBase) +(prezzoBase*uscite) +(prezzoBase/entrate) +(prezzoBase-uscite) = 0 \\ \hline
\end{tabular} \\
\end{center}
\begin{center}
\begin{tabular}{|p{11cm}|} \hline
\textbf{Output atteso}\\ \hline
\textless BusinessRule associated``Articolo'' comment``Comment'' name``BR'' rule``(entrate+prezzoBase)+ (prezzoBase*uscite)+ (prezzoBase/entrate)+ (prezzoBase-uscite)=0''\textgreater\\
\textless AstRuleVersion value``(= (+ (+ (+ (+ entrate prezzoBase) (* prezzoBase prezzoBase)) (/ prezzoBase entrate)) (- prezzoBase prezzoBase)) 0)''/\textgreater\\
\textless Bconf type````\textgreater\\
\textless OFloat type``+''\textgreater\\
\textless OFloat type``+''\textgreater\\
\textless OFloat type``+''\textgreater\\
\textless OFloat type``+''\textgreater\\
\textless FIELD\textgreater entrate\textless /FIELD\textgreater\\
\textless FIELD\textgreater prezzoBase\textless /FIELD\textgreater\\
\textless /OFloat\textgreater\\
\textless OFloat type``*''\textgreater\\
\textless FIELD\textgreater prezzoBase\textless /FIELD\textgreater\\
\textless FIELD\textgreater uscite\textless /FIELD\textgreater\\
\textless /OFloat\textgreater\\
\textless /OFloat\textgreater\\
\textless OFloat type``/''\textgreater\\
\textless FIELD\textgreater prezzoBase\textless /FIELD\textgreater\\
\textless FIELD\textgreater entrate\textless /FIELD\textgreater\\
\textless /OFloat\textgreater\\
\textless /OFloat\textgreater\\
\textless OFloat type``-''\textgreater\\
\textless FIELD\textgreater prezzoBase\textless /FIELD\textgreater\\
\textless FIELD\textgreater uscite\textless /FIELD\textgreater\\
\textless /OFloat\textgreater\\
\textless /OFloat\textgreater\\
\textless FLOAT\textgreater0\textless /FLOAT\textgreater\\
\textless /Bconf\textgreater\\
\textless /BusinessRule\textgreater\\
 \hline
\end{tabular} \\
\end{center}

\begin{Large}\textbf{VT7}\end{Large} \\
Si testa una \br\ per verificare che sia possibile fare la operazioni aritmetiche semplici tra vettori omogenei.\\
Requisiti soddisfatti dall'esecuzione: F1,F2,F10,NU1
\begin{center}
\begin{tabular}{|p{5cm}|p{6cm}|} \hline
\textbf{Attributo \br} & \textbf{Valore} \\ \hline
Nome & BR \\ \hline
Commento & Commento\\ \hline
Business object associato & Articolo \\ \hline
Regola & (prezzoBase+prezzoBase) +(prezzoBase*prezzoBase) +(prezzoBase/prezzoBase) +(prezzoBase-prezzoBase) =0 \\ \hline
\end{tabular} \\
\end{center}
\begin{center}
\begin{tabular}{|p{11cm}|} \hline
\textbf{Output atteso}\\ \hline
\textless BusinessRule associated``Articolo'' comment``Comment'' name``BR'' rule``(prezzoBase+prezzoBase)+ (prezzoBase*prezzoBase)+ (prezzoBase/prezzoBase)+ (prezzoBase-prezzoBase)=0''\textgreater\\
\textless AstRuleVersion value``(= (+ (+ (+ (+ prezzoBase prezzoBase) (* prezzoBase prezzoBase)) (/ prezzoBase prezzoBase)) (- prezzoBase prezzoBase)) 0)''/\textgreater\\
\textless Bconf type````\textgreater\\
\textless OFloat type``+''\textgreater\\
\textless OFloat type``+''\textgreater\\
\textless OFloat type``+''\textgreater\\
\textless OFloat type``+''\textgreater\\
\textless FIELD\textgreater prezzoBase\textless /FIELD\textgreater\\
\textless FIELD\textgreater prezzoBase\textless /FIELD\textgreater\\
\textless /OFloat\textgreater\\
\textless OFloat type``*''\textgreater\\
\textless FIELD\textgreater prezzoBase\textless /FIELD\textgreater\\
\textless FIELD\textgreater prezzoBase\textless /FIELD\textgreater\\
\textless /OFloat\textgreater\\
\textless /OFloat\textgreater\\
\textless OFloat type``/''\textgreater\\
\textless FIELD\textgreater prezzoBase\textless /FIELD\textgreater\\
\textless FIELD\textgreater prezzoBase\textless /FIELD\textgreater\\
\textless /OFloat\textgreater\\
\textless /OFloat\textgreater\\
\textless OFloat type``-''\textgreater\\
\textless FIELD\textgreater prezzoBase\textless /FIELD\textgreater\\
\textless FIELD\textgreater prezzoBase\textless /FIELD\textgreater\\
\textless /OFloat\textgreater\\
\textless /OFloat\textgreater\\
\textless FLOAT\textgreater0\textless /FLOAT\textgreater\\
\textless /Bconf\textgreater\\
\textless /BusinessRule\textgreater\\
 \hline
\end{tabular} \\
\end{center}

\begin{Large}\textbf{VT8a}\end{Large} \\
Si testa una \br\ per verificare che sia possibile fare la operazioni aritmetiche semplici tra matrici omogenee.\\
Requisiti soddisfatti dall'esecuzione: F1,F2,F10,NU1
\begin{center}
\begin{tabular}{|p{5cm}|p{6cm}|} \hline
\textbf{Attributo \br} & \textbf{Valore} \\ \hline
Nome & BR \\ \hline
Commento & Commento\\ \hline
Business object associato & Articolo \\ \hline
Regola & (matricePrezzi+matricePrezzi) +(matricePrezzi*matricePrezzi) +(matricePrezzi/matricePrezzi) +(matricePrezzi-matricePrezzi) =0 \\ \hline
\end{tabular} \\
\end{center}
\begin{center}
\begin{tabular}{|p{11cm}|} \hline
\textbf{Output atteso}\\ \hline
\textless BusinessRule associated``Articolo'' comment``Comment'' name``BR'' rule``(matricePrezzi+matricePrezzi)+ (matricePrezzi*matricePrezzi)+ (matricePrezzi/matricePrezzi)+ (matricePrezzi-matricePrezzi)=0''\textgreater\\
\textless AstRuleVersion value``(= (+ (+ (+ (+ matricePrezzi matricePrezzi) (* matricePrezzi matricePrezzi)) (/ matricePrezzi matricePrezzi)) (- matricePrezzi matricePrezzi)) 0)''/\textgreater\\
\textless Bconf type````\textgreater\\
\textless OFloat type``+''\textgreater\\
\textless OFloat type``+''\textgreater\\
\textless OFloat type``+''\textgreater\\
\textless OFloat type``+''\textgreater\\
\textless FIELD\textgreater matricePrezzi\textless /FIELD\textgreater\\
\textless FIELD\textgreater matricePrezzi\textless /FIELD\textgreater\\
\textless /OFloat\textgreater\\
\textless OFloat type``*''\textgreater\\
\textless FIELD\textgreater matricePrezzi\textless /FIELD\textgreater\\
\textless FIELD\textgreater matricePrezzi\textless /FIELD\textgreater\\
\textless /OFloat\textgreater\\
\textless /OFloat\textgreater\\
\textless OFloat type``/''\textgreater\\
\textless FIELD\textgreater matricePrezzi\textless /FIELD\textgreater\\
\textless FIELD\textgreater matricePrezzi\textless /FIELD\textgreater\\
\textless /OFloat\textgreater\\
\textless /OFloat\textgreater\\
\textless OFloat type``-''\textgreater\\
\textless FIELD\textgreater matricePrezzi\textless /FIELD\textgreater\\
\textless FIELD\textgreater matricePrezzi\textless /FIELD\textgreater\\
\textless /OFloat\textgreater\\
\textless /OFloat\textgreater\\
\textless FLOAT\textgreater0\textless /FLOAT\textgreater\\
\textless /Bconf\textgreater\\
\textless /BusinessRule\textgreater\\
 \hline
\end{tabular} \\
\end{center}

\begin{Large}\textbf{VT8b}\end{Large} \\
Si testa una \br\ per verificare che sia possibile fare la operazioni aritmetiche semplici tra matrici e scalari.\\
Requisiti soddisfatti dall'esecuzione: F1,F2,F10,NPr2,NQ1
\begin{center}
\begin{tabular}{|p{5cm}|p{6cm}|} \hline
\textbf{Attributo \br} & \textbf{Valore} \\ \hline
Nome & BR \\ \hline
Commento & Commento\\ \hline
Business object associato & Articolo \\ \hline
Regola & (matricePrezzi+2) +(matricePrezzi*2) +(matricePrezzi/2) +(matricePrezzi-2) =0 \\ \hline
\end{tabular} \\
\end{center}
\begin{center}
\begin{tabular}{|p{11cm}|} \hline
\textbf{Output atteso}\\ \hline
\textless BusinessRule associated``Articolo'' comment``Comment'' name``BR'' rule``(matricePrezzi+2)+ (matricePrezzi*2)+ (matricePrezzi/2)+ (matricePrezzi-2)=0''\textgreater\\
\textless AstRuleVersion value``(= (+ (+ (+ (+ matricePrezzi 2) (* matricePrezzi 2)) (/ matricePrezzi 2)) (- matricePrezzi 2)) 0)''/\textgreater\\
\textless Bconf type````\textgreater\\
\textless OFloat type``+''\textgreater\\
\textless OFloat type``+''\textgreater\\
\textless OFloat type``+''\textgreater\\
\textless OFloat type``+''\textgreater\\
\textless FIELD\textgreater matricePrezzi\textless /FIELD\textgreater\\
\textless FLOAT\textgreater 2\textless /FLOAT\textgreater\\
\textless /OFloat\textgreater\\
\textless OFloat type``*''\textgreater\\
\textless FIELD\textgreater matricePrezzi\textless /FIELD\textgreater\\
\textless FLOAT\textgreater 2\textless /FLOAT\textgreater\\
\textless /OFloat\textgreater\\
\textless /OFloat\textgreater\\
\textless OFloat type``/''\textgreater\\
\textless FIELD\textgreater matricePrezzi\textless /FIELD\textgreater\\
\textless FLOAT\textgreater 2\textless /FLOAT\textgreater\\
\textless /OFloat\textgreater\\
\textless /OFloat\textgreater\\
\textless OFloat type``-''\textgreater\\
\textless FIELD\textgreater matricePrezzi\textless /FIELD\textgreater\\
\textless FLOAT\textgreater 2\textless /FLOAT\textgreater\\
\textless /OFloat\textgreater\\
\textless /OFloat\textgreater\\
\textless FLOAT\textgreater0\textless /FLOAT\textgreater\\
\textless /Bconf\textgreater\\
\textless /BusinessRule\textgreater\\
 \hline
\end{tabular} \\
\end{center}

\begin{Large}\textbf{VT8c}\end{Large} \\
Si testa una \br\ per verificare che sia possibile fare la operazioni aritmetiche semplici tra matrici e campi dati scalari.\\
Requisiti soddisfatti dall'esecuzione: F1,F2,F10,NU1,NPr2,NQ1
\begin{center}
\begin{tabular}{|p{5cm}|p{6cm}|} \hline
\textbf{Attributo \br} & \textbf{Valore} \\ \hline
Nome & BR \\ \hline
Commento & Commento\\ \hline
Business object associato & Articolo \\ \hline
Regola & (matricePrezzi+entrate) +(matricePrezzi*entrate) +(matricePrezzi/entrate) +(matricePrezzi-entrate)=0 \\ \hline
\end{tabular} \\
\end{center}
\begin{center}
\begin{tabular}{|p{11cm}|} \hline
\textbf{Output atteso}\\ \hline
\textless BusinessRule associated``Articolo'' comment``Comment'' name``BR'' rule``(matricePrezzi+entrate)+ (matricePrezzi*entrate)+ (matricePrezzi/entrate)+ (matricePrezzi-entrate)=0''\textgreater\\
\textless AstRuleVersion value``(= (+ (+ (+ (+ matricePrezzi entrate) (* matricePrezzi entrate)) (/ matricePrezzi entrate)) (- matricePrezzi entrate)) 0)''/\textgreater\\
\textless Bconf type````\textgreater\\
\textless OFloat type``+''\textgreater\\
\textless OFloat type``+''\textgreater\\
\textless OFloat type``+''\textgreater\\
\textless OFloat type``+''\textgreater\\
\textless FIELD\textgreater matricePrezzi\textless /FIELD\textgreater\\
\textless FIELD\textgreater entrate\textless /FIELD\textgreater\\
\textless /OFloat\textgreater\\
\textless OFloat type``*''\textgreater\\
\textless FIELD\textgreater matricePrezzi\textless /FIELD\textgreater\\
\textless FIELD\textgreater entrate\textless /FIELD\textgreater\\
\textless /OFloat\textgreater\\
\textless /OFloat\textgreater\\
\textless OFloat type``/''\textgreater\\
\textless FIELD\textgreater matricePrezzi\textless /FIELD\textgreater\\
\textless FIELD\textgreater entrate\textless /FIELD\textgreater\\
\textless /OFloat\textgreater\\
\textless /OFloat\textgreater\\
\textless OFloat type``-''\textgreater\\
\textless FIELD\textgreater matricePrezzi\textless /FIELD\textgreater\\
\textless FIELD\textgreater entrate\textless /FIELD\textgreater\\
\textless /OFloat\textgreater\\
\textless /OFloat\textgreater\\
\textless FLOAT\textgreater0\textless /FLOAT\textgreater\\
\textless /Bconf\textgreater\\
\textless /BusinessRule\textgreater\\
 \hline
\end{tabular} \\
\end{center}

\begin{Large}\textbf{VT9a}\end{Large} \\
Si testa una \br\ per verificare che sia impossibile fare la operazioni aritmetiche semplici tra matrici e vettori.\\
Requisiti soddisfatti dall'esecuzione: F1,F2,F10,NU1
\begin{center}
\begin{tabular}{|p{5cm}|p{6cm}|} \hline
\textbf{Attributo \br} & \textbf{Valore} \\ \hline
Nome & BR \\ \hline
Commento & Commento\\ \hline
Business object associato & Articolo \\ \hline
Regola & (matricePrezzi+prezzoBase) +(matricePrezzi*prezzoBase) +(matricePrezzi/prezzoBase) +(matricePrezzi-prezzoBase) \\ \hline
\end{tabular} \\
\end{center}
\begin{center}
\begin{tabular}{|p{11cm}|} \hline
\textbf{Output atteso}\\ \hline
Operazione non consentita\\
 \hline
\end{tabular} \\
\end{center}

\begin{Large}\textbf{VT9b}\end{Large} \\
Si testa una \br\ per verificare che sia impossibile fare la operazioni aritmetiche semplici tra matrici non omogenee.\\
Requisiti soddisfatti dall'esecuzione:F1,F2,F10
\begin{center}
\begin{tabular}{|p{5cm}|p{6cm}|} \hline
\textbf{Attributo \br} & \textbf{Valore} \\ \hline
Nome & BR \\ \hline
Commento & Commento\\ \hline
Business object associato & Articolo \\ \hline
Regola & (matricePrezzi+grafico) +(matricePrezzi*grafico) +(matricePrezzi/grafico) +(matricePrezzi-grafico) \\ \hline
\end{tabular} \\
\end{center}
\begin{center}
\begin{tabular}{|p{11cm}|} \hline
\textbf{Output atteso}\\ \hline
Operazione non consentita\\
 \hline
\end{tabular} \\
\end{center}

\begin{Large}\textbf{VT9c}\end{Large} \\
Si testa una \br\ per verificare che sia impossibile fare la operazioni aritmetiche semplici tra matrici, vettori e scalari di tipo diverso.\\
Requisiti soddisfatti dall'esecuzione: F1,F2,F10
\begin{center}
\begin{tabular}{|p{5cm}|p{6cm}|} \hline
\textbf{Attributo \br} & \textbf{Valore} \\ \hline
Nome & BR \\ \hline
Commento & Commento\\ \hline
Business object associato & Articolo \\ \hline
Regola & matricePrezzi=``ArticoloUno'' AND prezziScontati=3 OR uscite=true \\ \hline
\end{tabular} \\
\end{center}
\begin{center}
\begin{tabular}{|p{11cm}|} \hline
\textbf{Output atteso}\\ \hline
Operazione non consentita\\
 \hline
\end{tabular} \\
\end{center}

\begin{Large}\textbf{VT10}\end{Large} \\
Si testa una \br\ per verificare che sia ben inteprtata la precedenza tra gli operatori.\\
Requisiti soddisfatti dall'esecuzione: F1,F2,F10
\begin{center}
\begin{tabular}{|p{5cm}|p{6cm}|} \hline
\textbf{Attributo \br} & \textbf{Valore} \\ \hline
Nome & BR \\ \hline
Commento & Commento\\ \hline
Business object associato & Articolo \\ \hline
Regola & uscite=2*3*4+4-3*5-5+3/2*3 \\ \hline
\end{tabular} \\
\end{center}
\begin{center}
\begin{tabular}{|p{11cm}|} \hline
\textbf{Output atteso}\\ \hline
\textless BusinessRule associated=``Articolo'' comment=``Commento'' name=``BR'' rule=``uscite=2*3*4+4-3*5-5+3/2*3''\textgreater\\
 \textless AstRuleVersion value=``(= uscite (+ (- (- (+ (* (* 2 3) 4) 4) (* 3 5)) 5) (* (/ 3 2) 3)))''/\textgreater\\
 \textless Bconf type=``=``\textgreater\\
 \textless FIELD\textgreater uscite\textless /FIELD\textgreater\\
 \textless OFloat type=``+''\textgreater\\
 \textless OFloat type=``-''\textgreater\\
 \textless OFloat type=``-''\textgreater \\
\textless OFloat type=``+''\textgreater \\
\textless OFloat type=``*''\textgreater\\
 \textless OFloat type=``*''\textgreater\\
 \textless FLOAT\textgreater 2\textless /FLOAT\textgreater\\
 \textless FLOAT\textgreater 3\textless /FLOAT\textgreater\\
 \textless /OFloat\textgreater\\
 \textless FLOAT\textgreater 4\textless /FLOAT\textgreater\\
 \textless /OFloat\textgreater\\
 \textless FLOAT\textgreater 4\textless /FLOAT\textgreater\\
 \textless /OFloat\textgreater\\
 \textless OFloat type=``*''\textgreater\\
 \textless FLOAT\textgreater 3\textless /FLOAT\textgreater\\
 \textless FLOAT\textgreater 5\textless /FLOAT\textgreater\\
 \textless /OFloat\textgreater\\
 \textless /OFloat\textgreater\\
 \textless FLOAT\textgreater 5\textless /FLOAT\textgreater\\
 \textless /OFloat\textgreater\\
 \textless OFloat type=``*''\textgreater\\
 \textless OFloat type=``/''\textgreater\\
 \textless FLOAT\textgreater 3\textless /FLOAT\textgreater\\
 \textless FLOAT\textgreater 2\textless /FLOAT\textgreater\\
 \textless /OFloat\textgreater\\
 \textless FLOAT\textgreater 3\textless /FLOAT\textgreater\\
 \textless /OFloat\textgreater\\
 \textless /OFloat\textgreater\\
 \textless /Bconf\textgreater\\
 \textless /BusinessRule\textgreater \\
 \hline
\end{tabular} \\
\end{center}

\begin{Large}\textbf{VT11}\end{Large} \\
Si testa una \br\ per verificare che siano ben intepretate le parentesi.\\
Requisiti soddisfatti dall'esecuzione: F1,F2,F10
\begin{center}
\begin{tabular}{|p{5cm}|p{6cm}|} \hline
\textbf{Attributo \br} & \textbf{Valore} \\ \hline
Nome & BR \\ \hline
Commento & Commento\\ \hline
Business object associato & Articolo \\ \hline
Regola & uscite=((((2*((3*4)+(4)-(3*5)-5)+3)/2)*3)+7) \\ \hline
\end{tabular} \\
\end{center}
\begin{center}
\begin{tabular}{|p{11cm}|} \hline
\textbf{Output atteso}\\ \hline
\textless BusinessRule associated=``Articolo'' comment=``Commento'' name=``BR'' rule=``uscite=((((2*((3*4)+(4)-(3*5)-5)+3)/2)*3)+7)''\textgreater\\
 \textless AstRuleVersion value=``(= uscite (+ (* (/ (+ (* 2 (- (- (+ (* 3 4) 4) (* 3 5)) 5)) 3) 2) 3) 7))''/\textgreater\\
 \textless Bconf type=``=``\textgreater \\
\textless FIELD\textgreater uscite\textless /FIELD\textgreater\\
 \textless OFloat type=``+''\textgreater\\
 \textless OFloat type=``*''\textgreater \\
\textless OFloat type=``/''\textgreater \\
\textless OFloat type=``+''\textgreater \\
\textless OFloat type=``*''\textgreater \\
\textless FLOAT\textgreater 2\textless /FLOAT\textgreater \\
\textless OFloat type=``-''\textgreater \\
\textless OFloat type=``-''\textgreater \\
\textless OFloat type=``+''\textgreater\\
 \textless OFloat type=``*''\textgreater \\
\textless FLOAT\textgreater 3\textless /FLOAT\textgreater \\
\textless FLOAT\textgreater 4\textless /FLOAT\textgreater \\
\textless /OFloat\textgreater \\
\textless FLOAT\textgreater 4\textless /FLOAT\textgreater \\
\textless /OFloat\textgreater \\
\textless OFloat type=``*''\textgreater \\
\textless FLOAT\textgreater 3\textless /FLOAT\textgreater \\
\textless FLOAT\textgreater 5\textless /FLOAT\textgreater \\
\textless /OFloat\textgreater \textless /OFloat\textgreater \\
\textless FLOAT\textgreater 5\textless /FLOAT\textgreater \\
\textless /OFloat\textgreater \textless /OFloat\textgreater \\
\textless FLOAT\textgreater 3\textless /FLOAT\textgreater \\
\textless /OFloat\textgreater \\
\textless FLOAT\textgreater 2\textless /FLOAT\textgreater \\
\textless /OFloat\textgreater \\
\textless FLOAT\textgreater 3\textless /FLOAT\textgreater\\
 \textless /OFloat\textgreater \\
\textless FLOAT\textgreater 7\textless /FLOAT\textgreater \\
\textless /OFloat\textgreater \\
\textless /Bconf\textgreater\\
 \textless /BusinessRule\textgreater \\
 \hline
\end{tabular} \\
\end{center}

\begin{Large}\textbf{VT12a}\end{Large} \\
Si testa una \br\ per verificare che sia impossibile effettuare operazioni tra le stringhe.\\
Requisiti soddisfatti dall'esecuzione: F1,F2,F10
\begin{center}
\begin{tabular}{|p{5cm}|p{6cm}|} \hline
\textbf{Attributo \br} & \textbf{Valore} \\ \hline
Nome & BR \\ \hline
Commento & Commento\\ \hline
Business object associato & Articolo \\ \hline
Regola & ``ArticoloUno''\textgreater``ArticoloUno'' AND ``ArticoloUno''\textgreater=``ArticoloUno'' AND ``ArticoloUno''\textless``ArticoloUno'' AND ``ArticoloUno''\textless=``ArticoloUno'' AND ``ArticoloUno''*nome  AND ``ArticoloUno''-nome  AND ``ArticoloUno''/nome  AND ``ArticoloUno''+nome\\ \hline
\end{tabular} \\
\end{center}
\begin{center}
\begin{tabular}{|p{11cm}|} \hline
\textbf{Output atteso}\\ \hline
Operazione non consentita\\
 \hline
\end{tabular} \\
\end{center}

\begin{Large}\textbf{VT12b}\end{Large} \\
Si testa una \br\ per verificare che sia possibile effettuare confronti semplici ('=', '!=') tra le stringhe.\\
Requisiti soddisfatti dall'esecuzione: F1,F2,F10
\begin{center}
\begin{tabular}{|p{5cm}|p{6cm}|} \hline
\textbf{Attributo \br} & \textbf{Valore} \\ \hline
Nome & BR \\ \hline
Commento & Commento\\ \hline
Business object associato & Articolo \\ \hline
Regola & nome=``ArticoloUno'' AND nome!=``ArticoloDue''\\ \hline
\end{tabular} \\
\end{center}
\begin{center}
\begin{tabular}{|p{11cm}|} \hline
\textbf{Output atteso}\\ \hline
\textless BusinessRule associated=``Articolo'' comment=``Commento'' name=``BR'' rule=``nome=\&quot;ArticoloUno\&quot;  AND nome!=\&quot;ArticoloDue\&quot;''\textgreater\\
 \textless AstRuleVersion value=``(AND (= nome \&quot;ArticoloUno\&quot;) (!= nome \&quot;ArticoloDue\&quot;))''/\textgreater\\
 \textless OpRule type=``AND''\textgreater\\
 \textless Conf type=``=``\textgreater\\
 \textless FIELD\textgreater nome\textless /FIELD\textgreater\\
 \textless STRING\textgreater ''ArticoloUno''\textless /STRING\textgreater\\
 \textless /Conf\textgreater\\
 \textless Conf type=``!=``\textgreater\\
 \textless FIELD\textgreater nome\textless /FIELD\textgreater\\
 \textless STRING\textgreater ''ArticoloDue''\textless /STRING\textgreater\\
 \textless /Conf\textgreater\\
 \textless /OpRule\textgreater\\
 \textless /BusinessRule\textgreater \\
 \hline
\end{tabular} \\
\end{center}

\begin{Large}\textbf{VT13a}\end{Large} \\
Si testa una \br\ per verificare che sia impossibile effettuare operazioni aritmetiche tra booleani.\\
Requisiti soddisfatti dall'esecuzione: F1,F2,F10
\begin{center}
\begin{tabular}{|p{5cm}|p{6cm}|} \hline
\textbf{Attributo \br} & \textbf{Valore} \\ \hline
Nome & BR \\ \hline
Commento & Commento\\ \hline
Business object associato & Articolo \\ \hline
Regola & (scontato+3) - (scontato*4)/scontato\\ \hline
\end{tabular} \\
\end{center}
\begin{center}
\begin{tabular}{|p{11cm}|} \hline
\textbf{Output atteso}\\ \hline
Operazione non consentita\\
 \hline
\end{tabular} \\
\end{center}

\begin{Large}\textbf{VT13b}\end{Large} \\
Si testa una \br\ per verificare che sia possibile effettuare operazioni logiche tra booleani.\\
Requisiti soddisfatti dall'esecuzione: F1,F2,F10
\begin{center}
\begin{tabular}{|p{5cm}|p{6cm}|} \hline
\textbf{Attributo \br} & \textbf{Valore} \\ \hline
Nome & BR \\ \hline
Commento & Commento\\ \hline
Business object associato & Articolo \\ \hline
Regola & ((scontato\textbar \textbar scontato)\&\& scontato) = true AND (true\textbar \textbar false\&\& NOT(true))=scontato\\ \hline
\end{tabular} \\
\end{center}
\begin{center}
\begin{tabular}{|p{11cm}|} \hline
\textbf{Output atteso}\\ \hline
\textless BusinessRule associated=``Articolo'' comment=``Commento'' name=``BR'' rule=``((scontato||scontato) \&amp;\&amp; scontato) = true AND (true || false \&amp;\&amp; NOT(true))=scontato''\textgreater \\
\textless AstRuleVersion value=``(AND (= (\&amp;\&amp; (|| scontato scontato) scontato) true) (= (\&amp;\&amp; (|| true false) (NOT true)) scontato))''/\textgreater\\
 \textless OpRule type=``AND''\textgreater\\
 \textless Conf type=``=``\textgreater\\
 \textless OBool type=``\&amp;\&amp;''\textgreater \\
\textless OBool type=``||''\textgreater\\
 \textless FIELD\textgreater scontato\textless /FIELD\textgreater\\
 \textless FIELD\textgreater scontato\textless /FIELD\textgreater\\
 \textless /OBool\textgreater\\
 \textless FIELD\textgreater scontato\textless /FIELD\textgreater\\
 \textless /OBool\textgreater\\
 \textless BOOL\textgreater true\textless /BOOL\textgreater\\
 \textless /Conf\textgreater\\
 \textless Conf type=``=``\textgreater\\
 \textless OBool type=``\&amp;\&amp;''\textgreater\\
 \textless OBool type=``||''\textgreater\\
 \textless BOOL\textgreater true\textless /BOOL\textgreater\\
 \textless BOOL\textgreater false\textless /BOOL\textgreater\\
 \textless /OBool\textgreater\\
 \textless BoFun type=``NOT''\textgreater\\
 \textless BOOL\textgreater true\textless /BOOL\textgreater\\
 \textless /BoFun\textgreater\\
 \textless /OBool\textgreater\\
 \textless FIELD\textgreater scontato\textless /FIELD\textgreater\\
 \textless /Conf\textgreater\\
 \textless /OpRule\textgreater\\
 \textless /BusinessRule\textgreater \\
 \hline
\end{tabular} \\
\end{center}

\begin{Large}\textbf{VT14}\end{Large} \\
Si testa una \br\ per verificare che sia possibile effettuare funzioni aritmetiche semplici come 'AVG', 'COUNT', 'SUM'.\\
Requisiti soddisfatti dall'esecuzione: F1,F2,F10,NU2
\begin{center}
\begin{tabular}{|p{5cm}|p{6cm}|} \hline
\textbf{Attributo \br} & \textbf{Valore} \\ \hline
Nome & BR \\ \hline
Commento & Commento\\ \hline
Business object associato & Articolo \\ \hline
Regola & SUM(prezzoBase) = 12 AND COUNT(entrate)=1, AND AVG(prezzoBase)=4\\ \hline
\end{tabular} \\
\end{center}
\begin{center}
\begin{tabular}{|p{11cm}|} \hline
\textbf{Output atteso}\\ \hline
\textless BusinessRule associated=``Articolo'' comment=``Commento'' name=``BR'' rule=``SUM(prezzoBase) = 12 AND COUNT(entrate)=1, AND AVG(prezzoBase)=4 ''\textgreater\\
 \textless AstRuleVersion value=``(AND (= (SUM prezzoBase) 12) (= (COUNT entrate) 1))''/\textgreater\\
 \textless OpRule type=``AND''\textgreater\\
 \textless Conf type=``=``\textgreater\\
 \textless FlFun type=``SUM''\textgreater\\
 \textless FIELD\textgreater prezzoBase\textless /FIELD\textgreater\\
 \textless /FlFun\textgreater\\
 \textless FLOAT\textgreater 12\textless /FLOAT\textgreater\\
 \textless /Conf\textgreater\\
 \textless Conf type=``=``\textgreater\\
 \textless Count type=``COUNT''\textgreater\\
 \textless FIELD\textgreater entrate\textless /FIELD\textgreater\\
 \textless /Count\textgreater\\
 \textless FLOAT\textgreater 1\textless /FLOAT\textgreater\\
 \textless /Conf\textgreater\\
 \textless /OpRule\textgreater\\
 \textless /BusinessRule\textgreater \\
 \hline
\end{tabular} \\
\end{center}

\begin{Large}\textbf{VT15a}\end{Large} \\
Si testa una \br\ per verificare che sia impossibile inserire regole prive di un operatore di confronto.\\
Requisiti soddisfatti dall'esecuzione: F1,F2,F10
\begin{center}
\begin{tabular}{|p{5cm}|p{6cm}|} \hline
\textbf{Attributo \br} & \textbf{Valore} \\ \hline
Nome & BR \\ \hline
Commento & Commento\\ \hline
Business object associato & Articolo \\ \hline
Regola & SUM(prezzoBase)\\ \hline
\end{tabular} \\
\end{center}
\begin{center}
\begin{tabular}{|p{11cm}|} \hline
\textbf{Output atteso}\\ \hline
Errore sintattico\\
 \hline
\end{tabular} \\
\end{center}

\begin{Large}\textbf{VT15b}\end{Large} \\
Si testa una \br\ per verificare che sia impossibile inserire regole prive di uno dei due elementi da confrontare\\
Requisiti soddisfatti dall'esecuzione: F1,F2,F10
\begin{center}
\begin{tabular}{|p{5cm}|p{6cm}|} \hline
\textbf{Attributo \br} & \textbf{Valore} \\ \hline
Nome & BR \\ \hline
Commento & Commento\\ \hline
Business object associato & Articolo \\ \hline
Regola & SUM(prezzoBase)=\\ \hline
\end{tabular} \\
\end{center}
\begin{center}
\begin{tabular}{|p{11cm}|} \hline
\textbf{Output atteso}\\ \hline
Errore sintattico\\
 \hline
\end{tabular} \\
\end{center}


\begin{Large}\textbf{VT15c}\end{Large} \\
Si testa una \br\ per verificare che sia impossibile inserire espressioni prive di operandi\\
Requisiti soddisfatti dall'esecuzione: F1,F2,F10
\begin{center}
\begin{tabular}{|p{5cm}|p{6cm}|} \hline
\textbf{Attributo \br} & \textbf{Valore} \\ \hline
Nome & BR \\ \hline
Commento & Commento\\ \hline
Business object associato & Articolo \\ \hline
Regola & 12 = 2entrate\\ \hline
\end{tabular} \\
\end{center}
\begin{center}
\begin{tabular}{|p{11cm}|} \hline
\textbf{Output atteso}\\ \hline
Errore sintattico\\
 \hline
\end{tabular} \\
\end{center}

\begin{Large}\textbf{VT15d}\end{Large} \\
Si testa una \br\ per verificare che sia impossibile inserire espressioni prive di operandi\\
Requisiti soddisfatti dall'esecuzione: F1,F2,F10
\begin{center}
\begin{tabular}{|p{5cm}|p{6cm}|} \hline
\textbf{Attributo \br} & \textbf{Valore} \\ \hline
Nome & BR \\ \hline
Commento & Commento\\ \hline
Business object associato & Articolo \\ \hline
Regola & uscite = 2 entrate\\ \hline
\end{tabular} \\
\end{center}
\begin{center}
\begin{tabular}{|p{11cm}|} \hline
\textbf{Output atteso}\\ \hline
Errore sintattico\\
 \hline
\end{tabular} \\
\end{center}

\begin{Large}\textbf{VT15e}\end{Large} \\
Si testa una \br\ per verificare che sia impossibile inserire regole non complete.\\
Requisiti soddisfatti dall'esecuzione: F1,F2,F10
\begin{center}
\begin{tabular}{|p{5cm}|p{6cm}|} \hline
\textbf{Attributo \br} & \textbf{Valore} \\ \hline
Nome & BR \\ \hline
Commento & Commento\\ \hline
Business object associato & Articolo \\ \hline
Regola & uscite = 2 AND\\ \hline
\end{tabular} \\
\end{center}
\begin{center}
\begin{tabular}{|p{11cm}|} \hline
\textbf{Output atteso}\\ \hline
Errore sintattico\\
 \hline
\end{tabular} \\
\end{center}

\begin{Large}\textbf{VT15f}\end{Large} \\
Si testa una \br\ per verificare che sia impossibile inserire parentesizzazioni errate.\\
Requisiti soddisfatti dall'esecuzione: F1,F2,F10
\begin{center}
\begin{tabular}{|p{5cm}|p{6cm}|} \hline
\textbf{Attributo \br} & \textbf{Valore} \\ \hline
Nome & BR \\ \hline
Commento & Commento\\ \hline
Business object associato & Articolo \\ \hline
Regola & ((2+4)=uscite\\ \hline
\end{tabular} \\
\end{center}
\begin{center}
\begin{tabular}{|p{11cm}|} \hline
\textbf{Output atteso}\\ \hline
Errore sintattico\\
 \hline
\end{tabular} \\
\end{center}

\begin{Large}\textbf{VT16}\end{Large} \\
Si testa una \br\ per verificare che siano supportato l'utilizzo dei numeri negativi.\\
Requisiti soddisfatti dall'esecuzione: F1,F2,F10
\begin{center}
\begin{tabular}{|p{5cm}|p{6cm}|} \hline
\textbf{Attributo \br} & \textbf{Valore} \\ \hline
Nome & BR \\ \hline
Commento & Commento\\ \hline
Business object associato & Articolo \\ \hline
Regola & uscite=((((\~2*((\~3*\~4)+(\~4)-(\~3*\~5)-5)+3)/\~2)*\~3)+7) \\ \hline
\end{tabular} \\
\end{center}
\begin{center}
\begin{tabular}{|p{11cm}|} \hline
\textbf{Output atteso}\\ \hline
\textless BusinessRule associated=``Articolo'' comment=``Commento'' name=``BR'' rule=``uscite=((((\~2*((\~3 *\~4)+(\~4)-(\~3*\~5) -5)+3)/\~2)*\~3)+7)''\textgreater \\
\textless AstRuleVersion value=``(= uscite (+ (* (/ (+ (* \~2 (- (- (+ (* \~3 \~4) \~4) (* \~3 \~5)) 5)) 3) \~2) \~3) 7))''/\textgreater\\ 
\textless Bconf type=``=``\textgreater \\
\textless FIELD\textgreater uscite\textless /FIELD\textgreater \\
\textless OFloat type=``+''\textgreater \\
\textless OFloat type=``*''\textgreater \\
\textless OFloat type=``/''\textgreater \\
\textless OFloat type=``+''\textgreater \\
\textless OFloat type=``*''\textgreater \\
\textless FLOAT\textgreater \~2\textless /FLOAT\textgreater \\
\textless OFloat type=``-''\textgreater \\
\textless OFloat type=``-''\textgreater \\
\textless OFloat type=``+''\textgreater \\
\textless OFloat type=``*''\textgreater \\
\textless FLOAT\textgreater \~3\textless /FLOAT\textgreater \\
\textless FLOAT\textgreater \~4\textless /FLOAT\textgreater \\
\textless /OFloat\textgreater \\
\textless FLOAT\textgreater \~4\textless /FLOAT\textgreater \\
\textless /OFloat\textgreater \\
\textless OFloat type=``*''\textgreater \\
\textless FLOAT\textgreater \~3\textless /FLOAT\textgreater \\
\textless FLOAT\textgreater \~5\textless /FLOAT\textgreater \\
\textless /OFloat\textgreater \\
\textless /OFloat\textgreater \\
\textless FLOAT\textgreater 5\textless /FLOAT\textgreater \\
\textless /OFloat\textgreater \\
\textless /OFloat\textgreater \\
\textless FLOAT\textgreater 3\textless /FLOAT\textgreater \\
\textless /OFloat\textgreater \\
\textless FLOAT\textgreater \~2\textless /FLOAT\textgreater \\
\textless /OFloat\textgreater \\
\textless FLOAT\textgreater \~3\textless /FLOAT\textgreater \\
\textless /OFloat\textgreater \\
\textless FLOAT\textgreater 7\textless /FLOAT\textgreater \\
\textless /OFloat\textgreater \\
\textless /Bconf\textgreater \\
\textless /BusinessRule\textgreater \\
 \hline
\end{tabular} \\
\end{center}

\begin{Large}\textbf{VT17}\end{Large} \\
Si testa una \br\ per verificare che siano ben tradotti tutti i caratteri non ammessi in un documento xml.\\
Requisiti soddisfatti dall'esecuzione: F1,F2,F10,NPo1
\begin{center}
\begin{tabular}{|p{5cm}|p{6cm}|} \hline
\textbf{Attributo \br} & \textbf{Valore} \\ \hline
Nome & \`a \`e \`i \`o \`u \`A \`E \`I \`O \`U ``test'' , 'test', \& , \textgreater , \textless \\ \hline
Commento & \`a \`e \`i \`o \`u \`A \`E \`I \`O \`U ``test'' , 'test', \& , \textgreater , \textless \\ \hline
Business object associato & Articolo \\ \hline
Regola & scontato=true \\ \hline
\end{tabular} \\
\end{center}
\begin{center}
\begin{tabular}{|p{11cm}|} \hline
\textbf{Output atteso}\\ \hline
\textless BusinessRule associated=``Articolo'' comment=``\&agrave; \&egrave; \&igrave; \&ograve; \&ugrave; \&Agrave; \&Egrave; \&Igrave; \&Ograve; \&Ugrave; \&quot;test\&quot;, 'test', \&amp;, \&gt;, \&lt;'' name=``\&agrave; \&egrave; \&igrave; \&ograve; \&ugrave; \&Agrave; \&Egrave; \&Igrave; \&Ograve; \&Ugrave; \&quot;test\&quot;, 'test', \&amp;, \&gt;, \&lt;'' rule=``scontato=true''\textgreater\\
 \textless AstRuleVersion value=``(= scontato true)''/\textgreater\\
\textless Bconf type=``=``\textgreater\\
 \textless FIELD\textgreater scontato\textless /FIELD\textgreater\\
 \textless BOOL\textgreater true\textless /BOOL\textgreater\\
 \textless /Bconf\textgreater\\
 \textless /BusinessRule\textgreater \\
 \hline
\end{tabular} \\
\end{center}

\begin{Large}\textbf{VT18}\end{Large} \\
Si testa una \br\ per verificare che siano ammessi solo alcuni caratteri nel message.\\
Requisiti soddisfatti dall'esecuzione: F1,F2,F10,NU3
\begin{center}
\begin{tabular}{|p{5cm}|p{6cm}|} \hline
\textbf{Attributo \br} & \textbf{Valore} \\ \hline
Nome & BR \\ \hline
Commento & Commento \\ \hline
Business object associato & Articolo \\ \hline
Regola & scontato=true message(``abc.. ABC.. 0..9 .?!\^\_-\~'') \\ \hline
\end{tabular} \\
\end{center}
\begin{center}
\begin{tabular}{|p{11cm}|} \hline
\textbf{Output atteso}\\ \hline
\textless BusinessRule associated=``Articolo'' comment=``Comment'' name=``BR'' rule=``scontato=true message(\&quote;abc.. ABC.. 0..9 .?!\&\#94; \&\#95;- \&\#126;\&quote;)''\textgreater\\
 \textless AstRuleVersion value=``(= scontato true) message(\&quote;abc.. ABC.. 0..9 .?!\&\#94; \&\#95;- \&\#126; \&quote;)''/\textgreater\\
\textless Bconf type=``=``\textgreater\\
 \textless FIELD\textgreater scontato\textless /FIELD\textgreater\\
 \textless BOOL\textgreater true\textless /BOOL\textgreater\\
 \textless /Bconf\textgreater\\
 \textless /BusinessRule\textgreater \\
 \hline
\end{tabular} \\
\end{center}

\begin{Large}\textbf{VT19}\end{Large} \\
Si testa una \br\ per verificare che non siano ammessi altri caratteri nel message.\\
Requisiti soddisfatti dall'esecuzione: F1,F2,F10
\begin{center}
\begin{tabular}{|p{5cm}|p{6cm}|} \hline
\textbf{Attributo \br} & \textbf{Valore} \\ \hline
Nome & BR \\ \hline
Commento & Commento \\ \hline
Business object associato & Articolo \\ \hline
Regola & scontato=true message(``\# []\{\}|\/=+-*\@\$\%'') \\ \hline
\end{tabular} \\
\end{center}
\begin{center}
\begin{tabular}{|p{11cm}|} \hline
\textbf{Output atteso}\\ \hline
Errore sintattico\\
 \hline
\end{tabular} \\
\end{center}

\begin{Large}\textbf{VT20}\end{Large} \\
Si testa una \br\ per verificare che non siano ammesse regole malformate.\\
Requisiti soddisfatti dall'esecuzione: F1,F2,F10
\begin{center}
\begin{tabular}{|p{5cm}|p{6cm}|} \hline
\textbf{Attributo \br} & \textbf{Valore} \\ \hline
Nome & BR \\ \hline
Commento & Commento \\ \hline
Business object associato & Articolo \\ \hline
Regola & scontato=true + \\ \hline
\end{tabular} \\
\end{center}
\begin{center}
\begin{tabular}{|p{11cm}|} \hline
\textbf{Output atteso}\\ \hline
Errore sintattico\\
 \hline
\end{tabular} \\
\end{center}

\subsection{Test sul comunicatore}
La macrocomponente comunicatore verr\`a testata a scatola aperta (whitebox), inserendo in input al comunicatore una stringa contenente una query \underline{XQuery} e verificando che la modifica sia stata effettuata nel repository. Non verranno testate le query non accettate dal linguaggio \underline{XQuery}, in quanto gia` garantite dal \underline{DBMS} \underline{eXist}. Per un corretto esito dei test \`e necessario che il \textit{repository} abbia una determinata situazione iniziale, illustrata nel riquadro sottostante. \\
\\

\begin{center}
\begin{tabular}{|p{11cm}|} \hline
\textless BusinessRule associated=``Articolo'' comment=``Commento'' name=``BR'' rule=``entrate=0 AND NOT(scontato)=true''\textgreater \\
\textless AstRuleVersion value=``(AND (= entrate 0) (= (NOT scontato) true))''/\textgreater \\
 \textless OpRule type=``AND''\textgreater \\
 \textless Conf type=``=''\textgreater \\
 \textless FIELD\textgreater entrate\textless /FIELD\textgreater \\
 \textless FLOAT\textgreater 0\textless /FLOAT\textgreater \\
 \textless /Conf\textgreater \\
 \textless Conf type=``=''\textgreater \\
\textless BoFun type=``NOT''\textgreater \\
 \textless FIELD\textgreater scontato\textless /FIELD\textgreater \\
\textless /BoFun\textgreater \\
 \textless BOOL\textgreater true\textless /BOOL\textgreater \\
\textless /Conf\textgreater \\
\textless /OpRule\textgreater \\
\textless /BusinessRule\textgreater \\ \hline
\textless BusinessRule associated=``Articolo'' comment=``Commento'' name=``BR1'' rule=``entrate=0 AND NOT(scontato)=true''\textgreater \\
\textless AstRuleVersion value=``(AND (= entrate 0) (= (NOT scontato) true))''/\textgreater \\
 \textless OpRule type=``AND''\textgreater \\
 \textless Conf type=``=''\textgreater \\
 \textless FIELD\textgreater entrate\textless /FIELD\textgreater \\
 \textless FLOAT\textgreater 0\textless /FLOAT\textgreater \\
 \textless /Conf\textgreater \\
 \textless Conf type=``=''\textgreater \\
\textless BoFun type=``NOT''\textgreater \\
 \textless FIELD\textgreater scontato\textless /FIELD\textgreater \\
\textless /BoFun\textgreater \\
 \textless BOOL\textgreater true\textless /BOOL\textgreater \\
\textless /Conf\textgreater \\
\textless /OpRule\textgreater \\
\textless /BusinessRule\textgreater \\ \hline
\end{tabular} \\
\end{center}

\begin{center}
\begin{tabular}{|p{11cm}|} \hline
\textless BusinessRule associated=``Articolo'' comment=``'' name=``BR2'' rule=``entrate=0 AND NOT(scontato)=true''\textgreater \\
\textless AstRuleVersion value=``(AND (= entrate 0) (= (NOT scontato) true))''/\textgreater \\
 \textless OpRule type=``AND''\textgreater \\
 \textless Conf type=``=''\textgreater \\
 \textless FIELD\textgreater entrate\textless /FIELD\textgreater \\
 \textless FLOAT\textgreater 0\textless /FLOAT\textgreater \\
 \textless /Conf\textgreater \\
 \textless Conf type=``=''\textgreater \\
\textless BoFun type=``NOT''\textgreater \\
 \textless FIELD\textgreater scontato\textless /FIELD\textgreater \\
\textless /BoFun\textgreater \\
 \textless BOOL\textgreater true\textless /BOOL\textgreater \\
\textless /Conf\textgreater \\
\textless /OpRule\textgreater \\
\textless /BusinessRule\textgreater \\ \hline
\textless BusinessRule associated=``Archivio'' comment=``'' name=``BR3'' rule=``entrate=0 AND NOT(scontato)=true''\textgreater \\
\textless AstRuleVersion value=``(AND (= entrate 0) (= (NOT scontato) true))''/\textgreater \\
 \textless OpRule type=``AND''\textgreater \\
 \textless Conf type=``=''\textgreater \\
 \textless FIELD\textgreater entrate\textless /FIELD\textgreater \\
 \textless FLOAT\textgreater 0\textless /FLOAT\textgreater \\
 \textless /Conf\textgreater \\
 \textless Conf type=``=''\textgreater \\
\textless BoFun type=``NOT''\textgreater \\
 \textless FIELD\textgreater scontato\textless /FIELD\textgreater \\
\textless /BoFun\textgreater \\
 \textless BOOL\textgreater true\textless /BOOL\textgreater \\
\textless /Conf\textgreater \\
\textless /OpRule\textgreater \\
\textless /BusinessRule\textgreater \\ \hline
\end{tabular} \\
\end{center}


\begin{Large}\textbf{CT1}\end{Large} \\
Si testa il comunicatore con una query che ritorni tutte le \br\ che hanno l'attributo name che comincia con BR.
Requisiti soddisfatti dall'esecuzione: F3,F5,F7,F8,NPo1,NPr2,NU3
\begin{center}
\begin{tabular}{|p{11cm}|} \hline
\textbf{Query eseguita}\\ \hline
 for \$br in //BusinessRule[starts-with(\@name,'BR')] return \$br\\ \hline
\end{tabular} \\
\end{center}
\begin{center}
\begin{tabular}{|p{11cm}|} \hline
\textbf{Output Atteso}\\ \hline
L'elenco delle \br\ in \underline{xml} come scritte nell'esempio di \rp\ sopra illustrato. In particolare dovranno essere visibili le \br\ : BR, BR1, BR2, BR3.\\ \hline
\end{tabular} \\
\end{center}

\begin{Large}\textbf{CT2}\end{Large} \\
Si testa il comunicatore con una query che cancelli tutte le \br\ il cui nome \`e BR1.
Requisiti soddisfatti dall'esecuzione: F3,F5,F7,F8,NPo1,NPr2,NU3
\begin{center}
\begin{tabular}{|p{11cm}|} \hline
\textbf{Query eseguita}\\ \hline
for \$br in //BusinessRule[\@name='BR1'] return update delete \$br;\\ \hline
\end{tabular} \\
\end{center}
\begin{center}
\begin{tabular}{|p{11cm}|} \hline
\textbf{Output Atteso}\\ \hline
La query BR1 deve essere rimossa dal \rp.\\ \hline
\end{tabular} \\
\end{center}

\begin{Large}\textbf{CT3}\end{Large} \\
Si testa il comunicatore con una query che ritorni il nome e la regola associata di tutte le \br\  che sono associate al \bo\ 'Articolo' e non hanno il commento.
Requisiti soddisfatti dall'esecuzione: F3,F5,F7,F8,NPo1,NPr2,NU3
\begin{center}
\begin{tabular}{|p{11cm}|} \hline
\textbf{Query eseguita}\\ \hline
for \$br in //BusinessRule[\@associated='Articolo' and not(\@comment)] return concat(\$br[\@name], \$br[\@rule])\\ \hline
\end{tabular} \\
\end{center}
\begin{center}
\begin{tabular}{|p{11cm}|} \hline
\textbf{Output Atteso}\\ \hline
Deve essere visibile la concatenazione del nome della regola seguito dalla regola stessa. In particolare ci\`o deve essere vero per le \br: BR, BR1.\\ \hline
\end{tabular} \\
\end{center}

\begin{Large}\textbf{CT4}\end{Large} \\
Si testa il comunicatore con una query che ritorni il nome e la regola associata di tutte le \br\  che non hanno il commento.
Requisiti soddisfatti dall'esecuzione: F3,F5,F7,F8,NPo1,NPr2,NU3
\begin{center}
\begin{tabular}{|p{11cm}|} \hline
\textbf{Query eseguita}\\ \hline
for \$br in //BusinessRule[not(\@comment)]return concat(\$br[\@name], \$br[\@rule])\\ \hline
\end{tabular} \\
\end{center}
\begin{center}
\begin{tabular}{|p{11cm}|} \hline
\textbf{Output Atteso}\\ \hline
Deve essere visibile la concatenazione del nome della regola seguito dalla regola stessa. In particolare ci\`o deve essere vero per le \br: BR2, BR3.\\ \hline
\end{tabular} \\
\end{center}

\begin{Large}\textbf{CT5}\end{Large} \\
Si testa il comunicatore con una query che inserisca una determinata \br .
Requisiti soddisfatti dall'esecuzione: F3,F5,F7,F8,NPo1,NPr2,NU3
\begin{center}
\begin{tabular}{|p{11cm}|} \hline
\textbf{Query eseguita}\\ \hline
update insert "\textless BusinessRule associated=``Archivio'' comment=``'' name=``BR3'' rule=``entrate=0 AND NOT(scontato)=true''\textgreater 
\textless AstRuleVersion value=``(AND (= entrate 0) (= (NOT scontato) true))''/\textgreater 
 \textless OpRule type=``AND''\textgreater 
 \textless Conf type=``=''\textgreater 
 \textless FIELD\textgreater entrate \textless /FIELD\textgreater 
 \textless FLOAT\textgreater 0 \textless /FLOAT\textgreater 
 \textless /Conf\textgreater 
 \textless Conf type=``='' \textgreater 
\textless BoFun type=``NOT'' \textgreater 
 \textless FIELD\textgreater scontato \textless /FIELD\textgreater 
\textless /BoFun\textgreater 
 \textless BOOL\textgreater true \textless /BOOL\textgreater 
\textless /Conf\textgreater 
\textless /OpRule\textgreater 
\textless /BusinessRule\textgreater " into /BusinessRules\\ \hline
\end{tabular} \\
\end{center}

\subsection{Test sulla GUI}
La macrocomponente gui verr\`a testata a scatola aperta (whitebox) e manualmente, verificando il funzionamento di ogni sua componente grafica. Verranno testati tutti i bottoni che lanciano nuove finestre. Si tester\`a che ogni bottone invochi il metodo appropiato.\\
\\
\begin{Large}\textbf{GT1a}\end{Large} \\
Si testa il bottone 'Inserisci Business Rule' e si testa che la finestra 'Inserisci' venga aperta.
Requisiti soddisfatti dall'esecuzione: F4,F8,NU3 \\
 \\
\begin{Large}\textbf{GT1b}\end{Large} \\
Si testa il bottone 'Rimuovi Business Rule' e si testa che la finestra 'Rimuovi' venga aperta.
Requisiti soddisfatti dall'esecuzione: F4,F8,NU3\\
 \\
\begin{Large}\textbf{GT1c}\end{Large} \\
Si testa il bottone 'Sandbox' e si testa che la finestra 'Sandbox' venga aperta.
Requisiti soddisfatti dall'esecuzione: F4,F8,NU3\\
 \\
\begin{Large}\textbf{GT2}\end{Large} \\
Si testa che tutti i campi textbox siano scrivibili.
Requisiti soddisfatti dall'esecuzione: F4,F8,NU3\\
 \\
\begin{Large}\textbf{GT3}\end{Large} \\
Si testa che ogni altro bottone della gui invochi il metodo corretto secondo le specifiche del sistema.
Requisiti soddisfatti dall'esecuzione: F4,F8,NU3\\
 \\
\section{Test di integrazione}
Vengono di seguito riportati i test riguardanti le integrazioni delle varie componenti. Questi test hanno come precondizione l'esecuzione dei test sulle singole macrocomponenti, di conseguenza sono mirati a scoprire difetti nell'interfacciamento e nell'interazione tra diverse componenti, nonch\`e di mostrare il soddisfacimento di alcuni requisiti.

\subsection{Integrazione validatore, comunicatore}
L'integrazione tra 'validatore' e 'comunicatore' verr\`a testata a scatola aperta (whitebox), inserendo in input al validatore un oggetto \textit{\br\ } con gli attributi di seguito specificati e verificando che l'output sia quello atteso. Essendo gi\`a state testate separatamente le 2 componenti ci si limiter\`a a testare le sole funzionalit\`a risultanti dalla loro integrazione.\\

\begin{Large}\textbf{GT3}\end{Large} \\
Si testa che ogni altro bottone della gui invochi il metodo corretto secondo le specifiche del sistema.
Requisiti soddisfatti dall'esecuzione: 

\subsection{Integrazione gui, validatore}
L'integrazione tra 'gui' e 'validatore' verr\`a testata a scatola aperta (whitebox), inserendo da gui  un oggetto \textit{\br\ } con gli attributi di seguito specificati e verificando che l'output sia quello atteso. Essendo gi\`a state testate separatamente le 2 componenti ci si limiter\`a a testare le sole funzionalit\`a risultanti dalla loro integrazione.\\

\subsection{Integrazione del sistema}
L'integrazionedell'intero sistema verr\`a testata a scatola aperta (whitebox), testando le sole funzionalit\`a risultanti dalla completa integrazione del sistema.\\

\section{Test di regressione}
Con test di regressione si intende la riescuzione di alcuni test case a fronte di una o pi\`u modifiche nel codice software. Ogni qual volta una componente venga modificata vanno rieseguiti (mediante l'utilizzo del sistema di automazione dei test) tutti i test riguardanti la componente specifica nonch\`e i test riguardanti l'integrazione della componente sessa con altre componenti fino all'esecuzione dei test di integrazione del sistema.
Negli esiti dei test, per non appesantire la lettura non vengono riportatitutti i test di regressione ma solo i test di regressione che abbiano prodotto un output o un esito differente.

\section{Tracciamento requisiti - test}
Nella seguente tabella vengono evidenziati i test necessari al soddisfacimento di ogni requisito. Gli identificativi dei test sono quelli riportati nel capitolo 'Descrizione dei test case'{link}. Mentre gli identificativi dei singoli requisiti sono quelli descritti nel documento \AR .
 \begin{center}
\begin{tabular}{|p{3cm}|p{8cm}|} \hline
\textbf{Requisito} & \textbf{Test che lo soddisfano}\\ \hline
F1  & VT1a, VT1b, VT1c, VT1d, VT1e, VT1f, VT2a, VT2b, VT2c, VT3a, VT3b, VT4, VT5a, VT5b, VT5c, VT6a, VT6b, VT7, VT8a, VT8b, VT8c, VT9a, VT9b, VT9c, VT10, VT11, VT12a, VT12b, VT13a, VT13b, VT14, VT15a, VT15b, VT15c, VT15d, VT15e, VT15f, VT16, VT17, VT18, VT19, VT20\\ \hline
F2  & VT1a, VT1b, VT1c, VT1d, VT1e, VT1f, VT2a, VT2b, VT2c, VT3a, VT3b, VT4, VT5a, VT5b, VT5c, VT6a, VT6b, VT7, VT8a, VT8b, VT8c, VT9a, VT9b, VT9c, VT10, VT11, VT12a, VT12b, VT13a, VT13b, VT14, VT15a, VT15b, VT15c, VT15d, VT15e, VT15f, VT16, VT17, VT18, VT19, VT20\\ \hline
F3  & CT1, CT2, CT3, CT4, CT5 \\ \hline
F4  & GT1a, GT1b, GT1c, GT2, GT3\\ \hline
F5  & CT1, CT2, CT3, CT4, CT5 \\ \hline
F6  &  \\ \hline
F7  & CT1, CT2, CT3, CT4, CT5  \\ \hline
F8  & CT1, CT2, CT3, CT4, CT5, GT1a, GT1b, GT1c, GT2, GT3  \\ \hline
F9  &  \\ \hline
F10 & VT1a,VT1b,VT1c,VT1d, VT2a, VT2b, VT2c, VT3a, VT3b, VT4, VT5a, VT5b, VT5c, VT6a, VT6b, VT7, VT8a, VT8b, VT8c, VT9a, VT9b, VT9c, VT10, VT11, VT12a, VT12b, VT13a, VT13b, VT14, VT15a, VT15b, VT15c, VT15d, VT15e, VT15f, VT16, VT17, VT18, VT19, VT20 \\ \hline
F11 & VT1a,VT1d \\ \hline
NU1 & VT6a, VT6b, VT7, VT8a, VT8c, VT9a\\ \hline
NU2 & VT14\\ \hline
NU3 & VT1a, VT3b, VT18, GT1a, GT1b, GT1c, GT2, GT3\\ \hline
NPo1 &  VT17, CT1, CT2, CT3, CT4, CT5 \\ \hline
NPr2 &  VT8b, VT8c, CT1, CT2, CT3, CT4, CT5 \\ \hline
NQ1 & VT8b, VT8c  \\ \hline
NQ2 &  \\ \hline
NQ3 &  \\ \hline
\end{tabular} \\
\end{center}


\chapter{Esiti dei test}
\section{Automatizzazione dei test}
Essendo la fase di test una fase costosa e laboriosa del processo software, si \`e deciso di automatizzare il pi\`e possibile l'esecuzione dei test. L'esecuzione dei test viene automatizzata mediante del semplice codice che si limita a invocare metodi e funzionalit\`a di una componente (o di alcune componenti integrate fra loro). Tale codice inizialmente si occupa di passare ai metodi delle componenti da testare degli opportuni parametri (alcuni errati ed altri corretti secondo gli input dei test case), confrontando poi l'output prodotto con quello atteso.
Alcuni test si rende necessario eseguirli a mano, come ad esempio i test riguardanti la componente \textit{gui}, tali test sono comunque una parte minima di tutti i test da eseguire non incidendo eccessivamente sul costo dell'intera fase di test.
\section{Esiti dei test dei difetti}
I \textit{test dei difetti} servono ad individuare possibili difetti o errori del codice su cui vengono testati. Per questo l'esito di un test dei difetti \`e segnato come negativo ogni qual volta non evidenzia un difetto, viene altres\`i segnato come positivo l'esito di quei test che abbiano individuato dei difetti nel software.
Vengono di seguito elencati gli esiti dei test effetuati. I test  sono presentati in forma tabellare, indicando il codice identificativo del test, seguito dall'esito del test. I test di regressione come illustrato precedentemente sono riportati solo dove abbiano prodotto un output o un esito differente dalla loro precedente esecuzione. I test di regressione sono evidenziati nelle tabelle sottostanti dalla sigla RT.


\subsection{Test sul validatore}
\begin{center}
\begin{tabular}{|p{1cm}|p{1.6cm}|p{8.4cm}|} \hline
\textbf{Id} & \textbf{Esito} & \textbf{Note} \\ \hline
VT1a & negativo & \\ \hline
VT1b & negativo & \\ \hline
VT1c & negativo & \\ \hline
VT1d & negativo & \\ \hline
VT1e & negativo & \\ \hline
VT1f & negativo & \\ \hline
VT2a & positivo & La parentesizzazione non \`e ancora stata introdotta nella grammatica del linguaggio. \\ \hline
VT2b & positivo & Il primo tag deve essere un 'OBool' e non un 'OpRule' .\\ \hline
VT2c & negativo & \\ \hline
VT3a & negativo & \\ \hline
VT3b & negativo & \\ \hline
VT4 & negativo & \\ \hline
VT5a & negativo & \\ \hline
VT5b & negativo & \\ \hline
VT5c & negativo & \\ \hline
VT6a & negativo & \\ \hline
VT6b & negativo & \\ \hline
VT7 & negativo & \\ \hline
VT8a & negativo & \\ \hline
VT8b & negativo & \\ \hline
VT8c & negativo & \\ \hline
VT9a & negativo & \\ \hline
VT9b & negativo & \\ \hline
VT9c & negativo & \\ \hline
VT10 & negativo & \\ \hline
VT11 & negativo & \\ \hline
VT12a & negativo & \\ \hline
VT12b & negativo & \\ \hline
VT13a & negativo & \\ \hline
VT13b & negativo & \\ \hline
VT14 & negativo & \\ \hline
\end{tabular} \\
\end{center}

\begin{center}
\begin{tabular}{|p{1cm}|p{1.6cm}|p{8.4cm}|} \hline
\textbf{Id} & \textbf{Esito} & \textbf{Note} \\ \hline
VT15a & negativo & \\ \hline
VT15b & negativo & \\ \hline
VT15c & positivo & \\ \hline
VT15d & positivo & \\ \hline
VT15e & negativo & \\ \hline
VT15f & negativo & \\ \hline
VT16 & negativo & \\ \hline
VT17 & positivo & Le lettere accentate non sono ancora state aggiunte nella grammatica del linguaggio.\\ \hline
VT18 & positivo & I numeri non sono ancora stati aggiunti alla grammatica del linguaggio.\\ \hline
VT19 & negativo & \\ \hline
VT20 & negativo & \\ \hline
\end{tabular} \\
\end{center}

\subsection{Test sul comunicatore}
\begin{center}
\begin{tabular}{|p{1cm}|p{1.6cm}|p{8.4cm}|} \hline
\textbf{Id} & \textbf{Esito} & \textbf{Note} \\ \hline
CT1 & negativo & \\ \hline
CT2 & negativo & \\ \hline
CT3 & negativo & \\ \hline
CT4 & negativo & \\ \hline
CT5 & negativo & \\ \hline
\end{tabular} \\
\end{center}

\subsection{Test sulla gui}
\begin{center}
\begin{tabular}{|p{1cm}|p{1.6cm}|p{8.4cm}|} \hline
\textbf{Id} & \textbf{Esito} & \textbf{Note} \\ \hline
GT1a & negativo & \\ \hline
GT1b & negativo & \\ \hline
GT1c & negativo & \\ \hline
GT2 & negativo & \\ \hline
GT3 & negativo & \\ \hline

\end{tabular} \\
\end{center}
\subsection{Integrazione validatore, comunicatore}
\subsection{Integrazione gui, validatore}
\subsection{Integrazione del sistema}

\chapter{Esiti dei test di convalida}
I \textit{test di convalida} servono ad dimostrare che tutti i requisiti sono stati soddisfatti. Per questo l'esito di un test di convalida \`e segnato come positivo ogni qual volta dimostra che un requisito \`e soddisfatto, viene altres\`i segnato come negativo l'esito di quei test che hanno fallito.
I test di convalida sono rimandati alla fase di verifica e validazione.
\end{document}
